% =================================================
% 3. AGENT SELECTION AND DEFINITION
% Included from main.tex with % =================================================
% 3. AGENT SELECTION AND DEFINITION
% Included from main.tex with % =================================================
% 3. AGENT SELECTION AND DEFINITION
% Included from main.tex with % =================================================
% 3. AGENT SELECTION AND DEFINITION
% Included from main.tex with \input{section3_agent_selection}
% (it replaces the empty Section 3 headings that are now in main.tex).
% Needs: tikz (arrows.meta, positioning), booktabs, tabularx, array, float -- already in
% the preamble. The diagram also needs shapes.geometric, loaded on the next line.
% =================================================
\usetikzlibrary{shapes.geometric}
\providecommand{\fname}[1]{\nolinkurl{#1}}

\section{Agent Selection and Definition}
\label{sec:agents}

The assignment fixes four crews (Data Management, Surveillance \& Conflict, Tactical
Resolution, and Reporting) and allows each team to add more crews and agents, as long as they are justified \cite{mas_assignment}. We designed the rest of the system together in the course forum and our messaging group.
Each proposal followed a common template (idea, motivation, agents and tools, place in the flow); the rest of the team agreed or answered with concerns, until we reached a solution that satisfied everyone. This section explains the result: which crews we
added and why (Section~\ref{sec:extra-crews}), what each crew does
(Section~\ref{sec:crews}), which agents form each crew (Section~\ref{sec:roles}) and how
everything is connected in the CrewAI Flow (Section~\ref{sec:overview}).
% -------------------------------------------------
\subsection{Decision on Additional Crews}
\label{sec:extra-crews}

Before discussing the proposals, we agreed on five rules to decide about them:

\begin{enumerate}
  \item \textbf{A new crew only for a new step in the flow.} If an idea needs its own input, its own output file and its own decision point, it becomes a crew. If it only adds information to a step that already exists, it becomes an agent inside that crew.
  \item \textbf{One detector.} Every conflict, between two aircraft or between an aircraft and a restricted zone, is found by the same deterministic tool. No crew checks separations with its own logic.
  \item \textbf{Numbers come from tools, decisions come from agents.} Distances, times and thresholds are always computed by deterministic tools. The language model interprets the results, prioritises and decides, but it never computes a separation.
  \item \textbf{Files as the interface.} Each crew reads the files of the previous crews and writes one intermediate file with a Pydantic model, so every crew can be developed and tested on its own, as the assignment recommends \cite{mas_assignment}.
  \item \textbf{Static inputs.} All the information is given as files at the start, and nothing is downloaded while the system runs.
\end{enumerate}

Table~\ref{tab:proposals} summarises the proposals we discussed and what we
decided about each one.

\begin{table}[!htbp]
\centering
\small
\caption{Design proposals and decisions.}
\label{tab:proposals}
\begin{tabularx}{\textwidth}{@{}>{\raggedright\arraybackslash}p{4.3cm}
                               >{\raggedright\arraybackslash}p{3.9cm} X@{}}
\toprule
\textbf{Proposal} & \textbf{Decision} & \textbf{Main reason} \\
\midrule
Meteorological hazards crew
  & New crew, extended to all airspace constraints
  & Wind changes the ETAs and storms create areas to avoid \\
Safety Validation crew
  & New crew, replaces the loop back to Surveillance
  & Tactical should not grade its own work \\
Restricted and military zones (Airspace Manager)
  & Agent in the Airspace Constraints crew
  & Not weather, but the same zone model \\
Airport constraints input
  & New input, handled by the Airspace Manager
  & A closure is a zone placed on a node \\
Flight Priority Officer
  & Agent in Data Management
  & It only adds the priority to the normalised data \\
Radio Communications Officer
  & Agent in Reporting
  & It only changes how the manoeuvres are written \\
Live weather download
  & Rejected: stored snapshot per scenario
  & Repeatable results and tests \\
Weather crew in parallel with Data Management
  & Rejected: sequential
  &  It needs the routes, levels and times from Data Management \\
Runway choice according to wind
  & Future work
  & Approach problem; our data has no runways \\
\bottomrule
\end{tabularx}
\end{table}

\textbf{Airspace Constraints crew.} This crew started as a proposal for meteorological
hazards. Wind changes the ground speed of each aircraft, so it changes the ETAs and
therefore when, and even whether, two aircraft meet. Thunderstorms also create areas that
aircraft must avoid: flight crews try to stay about 20~NM away from a cumulonimbus
\cite{skybrary_wx_avoidance}, and during weather avoidance speed control is not reliable for
separation \cite{skybrary_wx_los}. Later the team proposed two more kinds of restrictions:
restricted or military zones with active hours, which in Europe are managed with the
Flexible Use of Airspace concept \cite{eurocontrol_ernip3}, and airport constraints
(closures, ground stops, reduced capacity). At first we thought of adding the restricted
zones to the weather crew, but they have nothing to do with the weather. Instead of one
crew for each kind of restriction, we created one crew for all the constraints on the
airspace, with a specialised agent for each source. All of them produce the same object, a
\emph{constraint zone}, so the detector and the Safety Validation crew treat every
restriction in the same way.

This crew runs after Data Management, and not in parallel with it, because it needs its
output: the trajectories define the area, the flight levels and the time window that the
constraints must cover, and the first and last node of each route are the airports. It also
has to finish before Surveillance, because the detector needs the wind to compute correct ETAs
and the constraint zones to check the trajectories. The wind is used in two steps:
first the hazard zones, and then the wind-corrected ETAs, which need more changes in the
detector. Both steps are part of the design, but the second one will be implemented only
after the first one works.

\textbf{Safety Validation crew.} A manoeuvre can solve one conflict but create a new one,
or send an aircraft into a hazard zone. In real ATC a clearance is checked before it is
issued, and European controllers have short-term tools for conflict detection and
clearance verification \cite{skybrary_mtcd}. Our first draft of the flow sent the new
flight plans back to the Surveillance crew. We replaced this loop with a separate crew, so
that Tactical does not grade its own work and the feedback has a clear owner. To respect
rule~2, the Verification Officer uses the same tool as Surveillance and reads the
constraint zones from the file of the Airspace Constraints crew, without running that crew
again.

\textbf{Agents instead of crews.} Following rule~1, the other proposals did not need a new
step in the flow, so they became agents in existing crews: the Flight Priority Officer in
Data Management, the Airspace Manager in the Airspace Constraints crew and the Radio
Communications Officer in Reporting.

\textbf{Rejected ideas and future work.} We rejected downloading live weather data during
the execution: every run would give different results, it would be hard to test, and the
assignment gives the inputs as files. Instead, a small script downloads real data from
Open-Meteo \cite{openmeteo_docs} once and saves it in the scenario folder. Choosing the
runway according to the wind is left as future work, because our data has nodes but no
runways, and runway choice belongs to approach control, while our system works en-route
\cite{skybrary_area_controller}.

% -------------------------------------------------
\subsection{Crew Definitions}
\label{sec:crews}

The final system has six crews: the four mandatory ones and two additional ones (marked
with~*). Table~\ref{tab:crews} gives the purpose of each crew, its agents and the files it
reads and writes. All the intermediate files are written in the \fname{output/} folder.

\begin{table}[!htbp]
\centering
\small
\caption{Crews of the system. Files numbered \fname{01}--\fname{04} are intermediate
files in \fname{output/}.}
\label{tab:crews}
\begin{tabularx}{\textwidth}{@{}>{\raggedright\arraybackslash}p{2.3cm} X
                               >{\raggedright\arraybackslash}p{1.25cm}
                               >{\raggedright\arraybackslash}p{3.3cm}
                               >{\raggedright\arraybackslash}p{3.3cm}@{}}
\toprule
\textbf{Crew} & \textbf{Purpose} & \textbf{Agents} & \textbf{Reads} & \textbf{Writes} \\
\midrule
Data Management
  & Load, validate and normalise the inputs, build the nominal 4D trajectories and mark
    priority flights
  & FDA, TP, FPO
  & \fname{airspace.graphml}, \fname{aircraft.json}, \fname{flight_plans.json}
  & \fname{01_data_management.json} \\
\addlinespace
Airspace Constraints*
  & Turn the weather, the restricted zones and the airport constraints into wind per level
    and constraint zones
  & WDC, WHA, AM
  & \fname{01}, \fname{weather.json}, \fname{restricted_zones.json},
    \fname{airport_constraints.json}
  & \fname{01b_airspace_constraints.json} \\
\addlinespace
Surveillance \& Conflict
  & Detect every loss of separation (aircraft--aircraft and aircraft--zone) and prioritise
    the conflicts
  & RC, CA
  & \fname{01}, \fname{01b}
  & \fname{02_surveillance.json} \\
\addlinespace
Tactical Resolution
  & Propose one manoeuvre for each conflict
  & EC
  & \fname{01}, \fname{01b}, \fname{02}, \fname{04} (on a retry)
  & \fname{03_tactical_resolution.json} \\
\addlinespace
Safety Validation*
  & Check every manoeuvre again and approve or reject it with feedback
  & VO, SUP
  & \fname{01}, \fname{01b}, \fname{03}
  & \fname{04_safety_validation.json} \\
\addlinespace
Reporting
  & Write the ATC report and the radio instructions
  & RO, RCO
  & All previous files
  & \fname{report.md} \\
\bottomrule
\end{tabularx}
\end{table}

\textbf{Sequential processes with parallel tasks.} CrewAI also offers a hierarchical process,
where a manager agent delegates the tasks, as we saw in Lab~6 \cite{crewai_docs}. We do not use
it. Inside each crew the order of the tasks is always the same, and the decisions about what
happens next (which crew runs and how many retries are allowed) are taken by the routers of the
Flow, which are deterministic \cite{crewai_flows}. A manager agent would add one more
language-model decision exactly where we want the system to be predictable.

A sequential process does not mean that every task waits for the previous one, though.
Independent tasks run asynchronously, as we also saw in Lab~6. In the Airspace Constraints
crew, the Airspace Manager works in parallel with the weather agents, and in Data Management
the Trajectory Planner and the Flight Priority Officer both start as soon as the Flight Data
Assistant has validated the data. Their results are merged by deterministic code and not by
another agent, since all of them use the same Pydantic models.

This is also why Tactical Resolution has a single agent: one decision maker, checked by a
different crew. As future work, in larger scenarios the conflicts could be split into
independent groups, the connected components of the graph of conflicting aircraft
\cite{durand1995optimal, stollenwerk2017quantum}, and the Flow could run one Tactical Resolution
crew per group in parallel. If a manoeuvre created a new conflict between two groups, the
Safety Validation loop would merge them and solve them again. This
is similar to how real ATC divides the airspace into sectors with a manageable workload
\cite{skybrary_sectorisation}.

\textbf{Inputs.} Each scenario folder contains the three files required by the assignment and
three optional files that we added (Table~\ref{tab:inputs}). All times are given in seconds
from the start of the scenario. If an optional file is missing, the crew writes an empty list,
so the sample scenarios of the course still work. In cruise, aircraft fly at pressure altitude,
so we read the cruise altitude of each flight plan as a flight level (30000~ft = FL300), and
each flight level corresponds to a fixed pressure level in the weather data.

\begin{table}[!htbp]
\centering
\small
\caption{Input files of each scenario.}
\label{tab:inputs}
\begin{tabularx}{\textwidth}{@{}>{\raggedright\arraybackslash}p{3.6cm} X@{}}
\toprule
\textbf{File} & \textbf{Content} \\
\midrule
\multicolumn{2}{@{}l}{\textit{Required by the assignment}} \\
\fname{airspace.graphml}
  & Airspace graph. Each node is a waypoint with its latitude and longitude, and each edge
    is an airway between two waypoints with its distance in NM. \\
\addlinespace
\fname{aircraft.json}
   & One entry per aircraft: identifier and cruise speed in knots. We add an optional
    \fname{status} field with the special-handling indicator of the flight plan (for
    example \fname{STS/HOSP} or \fname{STS/MEDEVAC}) or the emergency transponder code
    \fname{7700} \cite{icao_doc4444}. \\
\addlinespace
\fname{flight_plans.json}
  & One entry per flight: aircraft identifier, route (ordered list of waypoints of the
    graph), departure time and cruise altitude in feet. The first and last waypoints of the
    route are its airports. \\
\midrule
\multicolumn{2}{@{}l}{\textit{Optional, added by us}} \\
\fname{weather.json}
  & Weather snapshot downloaded once from Open-Meteo \cite{openmeteo_docs}: the real date and
    time of the scenario start (\fname{scenario_start_utc}), and a grid of points that covers
    the airspace graph plus a 20~NM margin. For each point and hour, it gives the wind speed
    and direction at each pressure level near the cruise levels, the geopotential height of
    each level (its real altitude that day), the convective available potential energy
    (CAPE) and the freezing level. \\
\addlinespace
\fname{restricted_zones.json}
  & One entry per zone: identifier, type (restricted, military or danger area), shape (a
    circle with centre and radius in NM, or a polygon of latitude/longitude points), lower
    and upper flight level, and the time windows when it is active
    \cite{eurocontrol_ernip3}. \\
\addlinespace
\fname{airport_constraints.json}
  & One entry per constraint: airport node, type (closure, ground stop, reduced capacity,
    works or incident), affected movements (departures, arrivals or both), time window and,
    for reduced capacity, the maximum number of movements per hour. \\
\bottomrule
\end{tabularx}
\end{table}

The thresholds that turn the weather data into hazard zones are kept in one configuration
file, \fname{hazard_thresholds.yaml}, each with its reference: the CAPE above which an area
counts as a convective storm, the buffer around a storm (20~NM \cite{skybrary_wx_avoidance}),
the wind speed above which a level counts as turbulent, and the temperature range at which
icing is likely.

\textbf{Constraint zones.} All the restrictions share one model: an area (a circle or a
polygon, or a graph node for airports), an altitude band, a time window, a type and a
severity. Since the assignment does not define airports, we take the first and the last
node of each route as its airports. Airport constraints have no altitude band: a closure or
a ground stop affects the flights that depart from or arrive at that node inside the time
window, and Tactical solves it with a ground delay or a hold. Only the departure from the first node and the arrival at the last node of a flight plan
count as movements; flights that pass over the node at cruise level are not affected.

\textbf{Invalid flight plans.} If a flight plan refers to an unknown aircraft, a node that is
not in the graph or an airway that does not exist, FDA rejects it and gives the reason, as
the European flight plan processing system does \cite{eurocontrol_ifpl_spec}. The plan is not
repaired, because guessing a route would mean inventing data. The run continues with the
valid flights, and the rejected ones are excluded from detection and resolution, since they
have no trajectory. The report lists them with their reasons and warns that their conflicts
were not checked, because a rejected flight that still has to operate must have its plan
corrected and resubmitted \cite{eurocontrol_ifpl_spec}. If no flight plan is valid, the report
says that there was nothing to check instead of reporting that there are no conflicts.

\textbf{Weather coverage.} The weather snapshot must cover more than the planned routes,
because Tactical can reroute an aircraft, change its level or delay it. For this reason it
covers the whole airspace graph plus a 20~NM margin, every level that Tactical may assign
(cruise level $\pm$4000~ft) and the time window of the scenario including the maximum
allowed delay. In this way the Safety Validation crew never needs new weather data.

% -------------------------------------------------
\subsection{Agent Roles and Tools}
\label{sec:roles}

Each agent is inspired by a real role in air traffic management. As the assignment asks,
we keep the definitions short (role, goal and task) and avoid long backstories.
Table~\ref{tab:roles} defines the abbreviations of the agents used in the rest of the report.

\begin{table}[!htbp]
\centering
\small
\caption{Agents, the real role that inspires them, their main task and their tools.}
\label{tab:roles}
\begin{tabularx}{\textwidth}{@{}>{\raggedright\arraybackslash}p{2.6cm}
                               >{\raggedright\arraybackslash}p{2.75cm} X
                               >{\raggedright\arraybackslash}p{4.25cm}@{}}
\toprule
\textbf{Agent} & \textbf{Inspired by} & \textbf{Main task} & \textbf{Tools} \\
\midrule
\multicolumn{4}{@{}l}{\textit{Data Management crew}} \\
Flight Data Assistant (FDA)
  & Flight data assistant of an area control centre
  & Load the input files, validate the flight plans and normalise units and formats
  & \fname{ScenarioLoaderTool} \\
Trajectory Planner (TP)
  & Trajectory prediction of the flight data system \cite{skybrary_mtcd}
  & Build the nominal 4D trajectory of each flight (position and ETA at each waypoint)
  & \fname{TrajectoryTool} \\
Flight Priority Officer (FPO)
  & Supervisor handling emergency and priority flights \cite{icao_doc4444}
    & Turn the status codes into a priority level and handling rules (for example, do not
    move this aircraft), so that Tactical moves the other aircraft
  & \fname{PriorityRulesTool} \\
\midrule
\multicolumn{4}{@{}l}{\textit{Airspace Constraints crew}} \\
Weather Data Collector (WDC)
  & Meteorological office that supplies data to ATC
  & Load the weather snapshot and check that it covers the whole area, levels and time
  & \fname{WeatherSnapshotTool} \\
Weather Hazard Analyst (WHA)
  & Aviation forecaster issuing hazard warnings (SIGMET) \cite{aviationweather_api}
  & Apply the thresholds to the snapshot to get wind per flight level and hazard zones
    (storms, turbulence, icing), and describe them for the report
  & \fname{WeatherHazardTool} \\
Airspace Manager (AM)
  & Airspace Management Cell \cite{eurocontrol_ernip3}
  & Load restricted and military zones and airport constraints as constraint zones
  & \fname{AirspaceRestrictionTool} \\
\midrule
\multicolumn{4}{@{}l}{\textit{Surveillance \& Conflict crew}} \\
Radar Controller (RC)
  & Controller using conflict detection tools \cite{skybrary_mtcd}
  & Run the detector on all trajectories: aircraft against aircraft and against zones
  & \fname{ConflictDetectionTool} \\
Conflict Analyst (CA)
  & Planner controller \cite{skybrary_exec_planner}
  & Classify each conflict (type, severity, time left, priority flights) and order them
  & \fname{SeverityTool} \\
\midrule
\multicolumn{4}{@{}l}{\textit{Tactical Resolution crew}} \\
Executive Controller (EC)
  & Executive controller \cite{skybrary_exec_planner}
  & Choose one manoeuvre per conflict (level change, speed change, delay or hold,
    reroute), using the feedback of SUP on a retry
  & \fname{ManoeuvreSimulatorTool} \\
\midrule
\multicolumn{4}{@{}l}{\textit{Safety Validation crew}} \\
Verification Officer (VO)
  & Clearance verification \cite{skybrary_mtcd}
  & Run the detector again on the modified flight plans and the constraint zones
  & \fname{ConflictDetectionTool} \\
ATC Supervisor (SUP)
  & ATC shift supervisor
  & Approve or reject each manoeuvre and write structured feedback for EC
  & \fname{SeverityTool} \\
\midrule
\multicolumn{4}{@{}l}{\textit{Reporting crew}} \\
ATC Report Officer (RO)
  & Operational reporting of an ATC unit
  &  Write the Markdown report: situation, conflicts, resolutions, clearance status, flagged
    cases and rejected flight plans
  & \fname{FileReadTool} \\
Radio Communications Officer (RCO)
  & Radiotelephony of the executive controller \cite{skybrary_phraseology, icao_doc9432}
  & Turn each approved manoeuvre into a standard radio instruction
  & \fname{PhraseologyTool} \\
\bottomrule
\end{tabularx}
\end{table}

All the tools are deterministic. \fname{FileReadTool} is a built-in CrewAI tool, and the
other ten are custom tools built with \fname{BaseTool}, as in Lab~4 \cite{crewai_docs}.
\fname{ConflictDetectionTool} already exists: it wraps the detector of the course notebook
(Section~\ref{sec:agentification}) and will be extended to check constraint zones and wind.
Three tools deserve a short comment:

\begin{itemize}
  \item \fname{ManoeuvreSimulatorTool} is a ``what-if'' tool, like the one that supports
        MTCD \cite{skybrary_mtcd}. It applies a candidate manoeuvre to a flight plan,
        rebuilds the trajectory and calls the same detection functions, so EC can test an
        idea before proposing it. VO still checks the final proposal independently.
  \item The feedback of SUP is a Pydantic model with the conflict identifier, the decision,
        the reason (for example, ``new conflict with AC3 at 340~s'' or ``enters storm zone
        Z2'') and a suggestion, so that EC can actually use it on the retry.
  \item \fname{PhraseologyTool} fills fixed templates based on standard phraseology,
        because non-standard phraseology often contributes to incidents
        \cite{skybrary_phraseology}. The agent only chooses the template and the reason, for
        example ``AC2, climb to flight level 320, due traffic''.
\end{itemize}

EC can choose between four types of manoeuvre (Table~\ref{tab:manoeuvres}). Each one changes
the flight plan in a way that \fname{ManoeuvreSimulatorTool} can apply and check, and their
values use fixed steps, so the number of possible actions is finite.

\begin{table}[!htbp]
\centering
\small
\caption{Manoeuvres available to the Executive Controller.}
\label{tab:manoeuvres}
\begin{tabularx}{\textwidth}{@{}>{\raggedright\arraybackslash}p{2.4cm}
                               >{\raggedright\arraybackslash}p{4.6cm} X@{}}
\toprule
\textbf{Manoeuvre} & \textbf{What changes in the flight plan} & \textbf{Values and use} \\
\midrule
Level change
  & Cruise altitude
  & Climb or descend in steps of 1000~ft, within cruise level $\pm$4000~ft (the levels
    covered by the weather snapshot). Usually the first choice for traffic conflicts. \\
\addlinespace
Speed change
  & Speed of the aircraft
  & Steps of 0.01~Mach (about 6~kt at our cruise levels) \cite{skybrary_speed_control},
    at most $\pm$5\% of the planned speed, because at cruise levels the margin between the
    minimum and maximum speed is small \cite{skybrary_cruise_guide}. Not used for weather
    conflicts, because speed control is not reliable during weather avoidance
    \cite{skybrary_wx_los}. \\
\addlinespace
Delay or hold
  & Departure time (ground delay) or extra time at a waypoint (hold)
  & Whole minutes, up to the maximum delay covered by the snapshot. Used for airport
    constraints and when no level or route is free. \\
\addlinespace
Reroute
  & Route
  & Another path of the airspace graph between the same first and last node, so it only
    uses airways that exist. Used to avoid constraint zones. \\
\bottomrule
\end{tabularx}
\end{table}

% -------------------------------------------------
\subsection{System Overview Diagram}
\label{sec:overview}

Figure~\ref{fig:flow} shows the complete Flow: the six crews, the two routers and the
intermediate files.


The two routers are deterministic functions of the Flow \cite{crewai_flows}:

\begin{itemize}
  \item \textbf{Router~1} goes to Tactical Resolution if \fname{02_surveillance.json}
        contains at least one conflict, either between two aircraft or between an aircraft
        and a constraint zone. Otherwise it goes directly to Reporting, as the assignment
        requires \cite{mas_assignment}.
    \item \textbf{Router~2} goes to Reporting if every manoeuvre is approved. The iteration
        counter starts at 0 and counts how many times the flow has gone back to Tactical
        Resolution. If some manoeuvres are rejected and the counter is lower than two, the
        flow goes back to Tactical Resolution with the feedback of SUP, and only the rejected
        conflicts are solved again. If they are still rejected when the counter is two (the
        third attempt), the flow goes to Reporting and the conflict is flagged as ``requires
        human controller'', together with the last manoeuvre and the reason why it was
        rejected.
\end{itemize}

The iteration counter is stored in the state of the Flow and not in any agent, so no
language model can decide to skip the limit.

As an example, in the scenario \fname{sample_conflict} of the course, AC1 and AC2 fly in
opposite directions on the same route at 30000~ft. RC finds a head-on conflict at
$t = 112.5$~s and Router~1 sends it to Tactical. EC proposes that AC2 climbs to FL320. VO
runs the detector again and finds no conflict, because the vertical distance is now
2000~ft, so SUP approves it and Router~2 goes to Reporting. There, RCO writes the
instruction ``AC2, climb to flight level 320, due traffic'' and RO includes it in the
report with the clearance granted.

\begin{figure}[H]
\centering
\begin{tikzpicture}[
    font=\small,
    node distance=6.5mm and 10mm,
    crew/.style={draw, rounded corners=2pt, fill=gray!12, align=center,
                 minimum width=4.8cm, minimum height=9mm},
    extra/.style={crew, fill=blue!12},
    infile/.style={draw, align=center, font=\scriptsize, inner sep=3pt},
    router/.style={draw, diamond, aspect=2.4, fill=yellow!25, align=center,
                   inner sep=0.5pt, font=\scriptsize},
    arr/.style={-{Stealth[length=2mm]}, thick},
    lbl/.style={font=\scriptsize, inner sep=1.5pt}
  ]
  \node[infile, minimum width=4.8cm] (in)
    {\fname{airspace.graphml} \quad \fname{aircraft.json}\\ \fname{flight_plans.json}};
  \node[crew, below=of in] (dm) {Data Management Crew\\{\scriptsize FDA, TP, FPO}};
  \node[extra, below=9mm of dm] (ac) {Airspace Constraints Crew*\\{\scriptsize WDC, WHA, AM}};
  \node[infile, left=10mm of ac] (cin)
    {\fname{weather.json}\\ \fname{restricted_zones.json}\\ \fname{airport_constraints.json}};
  \node[crew, below=9mm of ac] (sc) {Surveillance \& Conflict Crew\\{\scriptsize RC, CA}};
  \node[router, below=of sc] (r1) {Router 1\\conflicts?};
  \node[crew, below=9mm of r1] (tr) {Tactical Resolution Crew\\{\scriptsize EC}};
  \node[extra, below=9mm of tr] (sv) {Safety Validation Crew*\\{\scriptsize VO, SUP}};
  \node[router, below=of sv] (r2) {Router 2\\approved?};
  \node[crew, below=12mm of r2] (rp) {Reporting Crew\\{\scriptsize RO, RCO}};
  \node[infile, below=of rp, double] (out) {\fname{report.md}};

  \draw[arr] (in) -- (dm);
  \draw[arr] (dm) -- node[lbl, right] {\fname{01_data_management.json}} (ac);
  \draw[arr] (cin) -- (ac);
  \draw[arr] (ac) -- node[lbl, right] {\fname{01b_airspace_constraints.json}} (sc);
  \draw[arr] (sc) -- node[lbl, right] {\fname{02_surveillance.json}} (r1);
  \draw[arr] (r1) -- node[lbl, right] {traffic or zone conflicts} (tr);
  \draw[arr] (tr) -- node[lbl, right] {\fname{03_tactical_resolution.json}} (sv);
  \draw[arr] (sv) -- node[lbl, right] {\fname{04_safety_validation.json}} (r2);
  \draw[arr] (r2) -- node[lbl, right, align=left]
    {all approved, or still rejected\\at iteration 2 (flagged)} (rp);
  \draw[arr] (rp) -- (out);

  \draw[arr] (r1.east) -- ++(4.2,0) coordinate (c1)
    |- node[lbl, pos=0.25, right] {no conflicts} (rp.east);
  \draw[arr] (r2.west) -- ++(-2.9,0)
    |- node[lbl, pos=0.25, left, align=right] {rejected,\\iteration $<2$} (tr.west);
\end{tikzpicture}
\caption{CrewAI Flow of the system. Blue crews (*) are the two additional crews. Every crew
can also read the earlier intermediate files; for example, Tactical Resolution and Safety
Validation read the constraint zones of \fname{01b}.}
\label{fig:flow}
\end{figure}

% (it replaces the empty Section 3 headings that are now in main.tex).
% Needs: tikz (arrows.meta, positioning), booktabs, tabularx, array, float -- already in
% the preamble. The diagram also needs shapes.geometric, loaded on the next line.
% =================================================
\usetikzlibrary{shapes.geometric}
\providecommand{\fname}[1]{\nolinkurl{#1}}

\section{Agent Selection and Definition}
\label{sec:agents}

The assignment fixes four crews (Data Management, Surveillance \& Conflict, Tactical
Resolution, and Reporting) and allows each team to add more crews and agents, as long as they are justified \cite{mas_assignment}. We designed the rest of the system together in the course forum and our messaging group.
Each proposal followed a common template (idea, motivation, agents and tools, place in the flow); the rest of the team agreed or answered with concerns, until we reached a solution that satisfied everyone. This section explains the result: which crews we
added and why (Section~\ref{sec:extra-crews}), what each crew does
(Section~\ref{sec:crews}), which agents form each crew (Section~\ref{sec:roles}) and how
everything is connected in the CrewAI Flow (Section~\ref{sec:overview}).
% -------------------------------------------------
\subsection{Decision on Additional Crews}
\label{sec:extra-crews}

Before discussing the proposals, we agreed on five rules to decide about them:

\begin{enumerate}
  \item \textbf{A new crew only for a new step in the flow.} If an idea needs its own input, its own output file and its own decision point, it becomes a crew. If it only adds information to a step that already exists, it becomes an agent inside that crew.
  \item \textbf{One detector.} Every conflict, between two aircraft or between an aircraft and a restricted zone, is found by the same deterministic tool. No crew checks separations with its own logic.
  \item \textbf{Numbers come from tools, decisions come from agents.} Distances, times and thresholds are always computed by deterministic tools. The language model interprets the results, prioritises and decides, but it never computes a separation.
  \item \textbf{Files as the interface.} Each crew reads the files of the previous crews and writes one intermediate file with a Pydantic model, so every crew can be developed and tested on its own, as the assignment recommends \cite{mas_assignment}.
  \item \textbf{Static inputs.} All the information is given as files at the start, and nothing is downloaded while the system runs.
\end{enumerate}

Table~\ref{tab:proposals} summarises the proposals we discussed and what we
decided about each one.

\begin{table}[!htbp]
\centering
\small
\caption{Design proposals and decisions.}
\label{tab:proposals}
\begin{tabularx}{\textwidth}{@{}>{\raggedright\arraybackslash}p{4.3cm}
                               >{\raggedright\arraybackslash}p{3.9cm} X@{}}
\toprule
\textbf{Proposal} & \textbf{Decision} & \textbf{Main reason} \\
\midrule
Meteorological hazards crew
  & New crew, extended to all airspace constraints
  & Wind changes the ETAs and storms create areas to avoid \\
Safety Validation crew
  & New crew, replaces the loop back to Surveillance
  & Tactical should not grade its own work \\
Restricted and military zones (Airspace Manager)
  & Agent in the Airspace Constraints crew
  & Not weather, but the same zone model \\
Airport constraints input
  & New input, handled by the Airspace Manager
  & A closure is a zone placed on a node \\
Flight Priority Officer
  & Agent in Data Management
  & It only adds the priority to the normalised data \\
Radio Communications Officer
  & Agent in Reporting
  & It only changes how the manoeuvres are written \\
Live weather download
  & Rejected: stored snapshot per scenario
  & Repeatable results and tests \\
Weather crew in parallel with Data Management
  & Rejected: sequential
  &  It needs the routes, levels and times from Data Management \\
Runway choice according to wind
  & Future work
  & Approach problem; our data has no runways \\
\bottomrule
\end{tabularx}
\end{table}

\textbf{Airspace Constraints crew.} This crew started as a proposal for meteorological
hazards. Wind changes the ground speed of each aircraft, so it changes the ETAs and
therefore when, and even whether, two aircraft meet. Thunderstorms also create areas that
aircraft must avoid: flight crews try to stay about 20~NM away from a cumulonimbus
\cite{skybrary_wx_avoidance}, and during weather avoidance speed control is not reliable for
separation \cite{skybrary_wx_los}. Later the team proposed two more kinds of restrictions:
restricted or military zones with active hours, which in Europe are managed with the
Flexible Use of Airspace concept \cite{eurocontrol_ernip3}, and airport constraints
(closures, ground stops, reduced capacity). At first we thought of adding the restricted
zones to the weather crew, but they have nothing to do with the weather. Instead of one
crew for each kind of restriction, we created one crew for all the constraints on the
airspace, with a specialised agent for each source. All of them produce the same object, a
\emph{constraint zone}, so the detector and the Safety Validation crew treat every
restriction in the same way.

This crew runs after Data Management, and not in parallel with it, because it needs its
output: the trajectories define the area, the flight levels and the time window that the
constraints must cover, and the first and last node of each route are the airports. It also
has to finish before Surveillance, because the detector needs the wind to compute correct ETAs
and the constraint zones to check the trajectories. The wind is used in two steps:
first the hazard zones, and then the wind-corrected ETAs, which need more changes in the
detector. Both steps are part of the design, but the second one will be implemented only
after the first one works.

\textbf{Safety Validation crew.} A manoeuvre can solve one conflict but create a new one,
or send an aircraft into a hazard zone. In real ATC a clearance is checked before it is
issued, and European controllers have short-term tools for conflict detection and
clearance verification \cite{skybrary_mtcd}. Our first draft of the flow sent the new
flight plans back to the Surveillance crew. We replaced this loop with a separate crew, so
that Tactical does not grade its own work and the feedback has a clear owner. To respect
rule~2, the Verification Officer uses the same tool as Surveillance and reads the
constraint zones from the file of the Airspace Constraints crew, without running that crew
again.

\textbf{Agents instead of crews.} Following rule~1, the other proposals did not need a new
step in the flow, so they became agents in existing crews: the Flight Priority Officer in
Data Management, the Airspace Manager in the Airspace Constraints crew and the Radio
Communications Officer in Reporting.

\textbf{Rejected ideas and future work.} We rejected downloading live weather data during
the execution: every run would give different results, it would be hard to test, and the
assignment gives the inputs as files. Instead, a small script downloads real data from
Open-Meteo \cite{openmeteo_docs} once and saves it in the scenario folder. Choosing the
runway according to the wind is left as future work, because our data has nodes but no
runways, and runway choice belongs to approach control, while our system works en-route
\cite{skybrary_area_controller}.

% -------------------------------------------------
\subsection{Crew Definitions}
\label{sec:crews}

The final system has six crews: the four mandatory ones and two additional ones (marked
with~*). Table~\ref{tab:crews} gives the purpose of each crew, its agents and the files it
reads and writes. All the intermediate files are written in the \fname{output/} folder.

\begin{table}[!htbp]
\centering
\small
\caption{Crews of the system. Files numbered \fname{01}--\fname{04} are intermediate
files in \fname{output/}.}
\label{tab:crews}
\begin{tabularx}{\textwidth}{@{}>{\raggedright\arraybackslash}p{2.3cm} X
                               >{\raggedright\arraybackslash}p{1.25cm}
                               >{\raggedright\arraybackslash}p{3.3cm}
                               >{\raggedright\arraybackslash}p{3.3cm}@{}}
\toprule
\textbf{Crew} & \textbf{Purpose} & \textbf{Agents} & \textbf{Reads} & \textbf{Writes} \\
\midrule
Data Management
  & Load, validate and normalise the inputs, build the nominal 4D trajectories and mark
    priority flights
  & FDA, TP, FPO
  & \fname{airspace.graphml}, \fname{aircraft.json}, \fname{flight_plans.json}
  & \fname{01_data_management.json} \\
\addlinespace
Airspace Constraints*
  & Turn the weather, the restricted zones and the airport constraints into wind per level
    and constraint zones
  & WDC, WHA, AM
  & \fname{01}, \fname{weather.json}, \fname{restricted_zones.json},
    \fname{airport_constraints.json}
  & \fname{01b_airspace_constraints.json} \\
\addlinespace
Surveillance \& Conflict
  & Detect every loss of separation (aircraft--aircraft and aircraft--zone) and prioritise
    the conflicts
  & RC, CA
  & \fname{01}, \fname{01b}
  & \fname{02_surveillance.json} \\
\addlinespace
Tactical Resolution
  & Propose one manoeuvre for each conflict
  & EC
  & \fname{01}, \fname{01b}, \fname{02}, \fname{04} (on a retry)
  & \fname{03_tactical_resolution.json} \\
\addlinespace
Safety Validation*
  & Check every manoeuvre again and approve or reject it with feedback
  & VO, SUP
  & \fname{01}, \fname{01b}, \fname{03}
  & \fname{04_safety_validation.json} \\
\addlinespace
Reporting
  & Write the ATC report and the radio instructions
  & RO, RCO
  & All previous files
  & \fname{report.md} \\
\bottomrule
\end{tabularx}
\end{table}

\textbf{Sequential processes with parallel tasks.} CrewAI also offers a hierarchical process,
where a manager agent delegates the tasks, as we saw in Lab~6 \cite{crewai_docs}. We do not use
it. Inside each crew the order of the tasks is always the same, and the decisions about what
happens next (which crew runs and how many retries are allowed) are taken by the routers of the
Flow, which are deterministic \cite{crewai_flows}. A manager agent would add one more
language-model decision exactly where we want the system to be predictable.

A sequential process does not mean that every task waits for the previous one, though.
Independent tasks run asynchronously, as we also saw in Lab~6. In the Airspace Constraints
crew, the Airspace Manager works in parallel with the weather agents, and in Data Management
the Trajectory Planner and the Flight Priority Officer both start as soon as the Flight Data
Assistant has validated the data. Their results are merged by deterministic code and not by
another agent, since all of them use the same Pydantic models.

This is also why Tactical Resolution has a single agent: one decision maker, checked by a
different crew. As future work, in larger scenarios the conflicts could be split into
independent groups, the connected components of the graph of conflicting aircraft
\cite{durand1995optimal, stollenwerk2017quantum}, and the Flow could run one Tactical Resolution
crew per group in parallel. If a manoeuvre created a new conflict between two groups, the
Safety Validation loop would merge them and solve them again. This
is similar to how real ATC divides the airspace into sectors with a manageable workload
\cite{skybrary_sectorisation}.

\textbf{Inputs.} Each scenario folder contains the three files required by the assignment and
three optional files that we added (Table~\ref{tab:inputs}). All times are given in seconds
from the start of the scenario. If an optional file is missing, the crew writes an empty list,
so the sample scenarios of the course still work. In cruise, aircraft fly at pressure altitude,
so we read the cruise altitude of each flight plan as a flight level (30000~ft = FL300), and
each flight level corresponds to a fixed pressure level in the weather data.

\begin{table}[!htbp]
\centering
\small
\caption{Input files of each scenario.}
\label{tab:inputs}
\begin{tabularx}{\textwidth}{@{}>{\raggedright\arraybackslash}p{3.6cm} X@{}}
\toprule
\textbf{File} & \textbf{Content} \\
\midrule
\multicolumn{2}{@{}l}{\textit{Required by the assignment}} \\
\fname{airspace.graphml}
  & Airspace graph. Each node is a waypoint with its latitude and longitude, and each edge
    is an airway between two waypoints with its distance in NM. \\
\addlinespace
\fname{aircraft.json}
   & One entry per aircraft: identifier and cruise speed in knots. We add an optional
    \fname{status} field with the special-handling indicator of the flight plan (for
    example \fname{STS/HOSP} or \fname{STS/MEDEVAC}) or the emergency transponder code
    \fname{7700} \cite{icao_doc4444}. \\
\addlinespace
\fname{flight_plans.json}
  & One entry per flight: aircraft identifier, route (ordered list of waypoints of the
    graph), departure time and cruise altitude in feet. The first and last waypoints of the
    route are its airports. \\
\midrule
\multicolumn{2}{@{}l}{\textit{Optional, added by us}} \\
\fname{weather.json}
  & Weather snapshot downloaded once from Open-Meteo \cite{openmeteo_docs}: the real date and
    time of the scenario start (\fname{scenario_start_utc}), and a grid of points that covers
    the airspace graph plus a 20~NM margin. For each point and hour, it gives the wind speed
    and direction at each pressure level near the cruise levels, the geopotential height of
    each level (its real altitude that day), the convective available potential energy
    (CAPE) and the freezing level. \\
\addlinespace
\fname{restricted_zones.json}
  & One entry per zone: identifier, type (restricted, military or danger area), shape (a
    circle with centre and radius in NM, or a polygon of latitude/longitude points), lower
    and upper flight level, and the time windows when it is active
    \cite{eurocontrol_ernip3}. \\
\addlinespace
\fname{airport_constraints.json}
  & One entry per constraint: airport node, type (closure, ground stop, reduced capacity,
    works or incident), affected movements (departures, arrivals or both), time window and,
    for reduced capacity, the maximum number of movements per hour. \\
\bottomrule
\end{tabularx}
\end{table}

The thresholds that turn the weather data into hazard zones are kept in one configuration
file, \fname{hazard_thresholds.yaml}, each with its reference: the CAPE above which an area
counts as a convective storm, the buffer around a storm (20~NM \cite{skybrary_wx_avoidance}),
the wind speed above which a level counts as turbulent, and the temperature range at which
icing is likely.

\textbf{Constraint zones.} All the restrictions share one model: an area (a circle or a
polygon, or a graph node for airports), an altitude band, a time window, a type and a
severity. Since the assignment does not define airports, we take the first and the last
node of each route as its airports. Airport constraints have no altitude band: a closure or
a ground stop affects the flights that depart from or arrive at that node inside the time
window, and Tactical solves it with a ground delay or a hold. Only the departure from the first node and the arrival at the last node of a flight plan
count as movements; flights that pass over the node at cruise level are not affected.

\textbf{Invalid flight plans.} If a flight plan refers to an unknown aircraft, a node that is
not in the graph or an airway that does not exist, FDA rejects it and gives the reason, as
the European flight plan processing system does \cite{eurocontrol_ifpl_spec}. The plan is not
repaired, because guessing a route would mean inventing data. The run continues with the
valid flights, and the rejected ones are excluded from detection and resolution, since they
have no trajectory. The report lists them with their reasons and warns that their conflicts
were not checked, because a rejected flight that still has to operate must have its plan
corrected and resubmitted \cite{eurocontrol_ifpl_spec}. If no flight plan is valid, the report
says that there was nothing to check instead of reporting that there are no conflicts.

\textbf{Weather coverage.} The weather snapshot must cover more than the planned routes,
because Tactical can reroute an aircraft, change its level or delay it. For this reason it
covers the whole airspace graph plus a 20~NM margin, every level that Tactical may assign
(cruise level $\pm$4000~ft) and the time window of the scenario including the maximum
allowed delay. In this way the Safety Validation crew never needs new weather data.

% -------------------------------------------------
\subsection{Agent Roles and Tools}
\label{sec:roles}

Each agent is inspired by a real role in air traffic management. As the assignment asks,
we keep the definitions short (role, goal and task) and avoid long backstories.
Table~\ref{tab:roles} defines the abbreviations of the agents used in the rest of the report.

\begin{table}[!htbp]
\centering
\small
\caption{Agents, the real role that inspires them, their main task and their tools.}
\label{tab:roles}
\begin{tabularx}{\textwidth}{@{}>{\raggedright\arraybackslash}p{2.6cm}
                               >{\raggedright\arraybackslash}p{2.75cm} X
                               >{\raggedright\arraybackslash}p{4.25cm}@{}}
\toprule
\textbf{Agent} & \textbf{Inspired by} & \textbf{Main task} & \textbf{Tools} \\
\midrule
\multicolumn{4}{@{}l}{\textit{Data Management crew}} \\
Flight Data Assistant (FDA)
  & Flight data assistant of an area control centre
  & Load the input files, validate the flight plans and normalise units and formats
  & \fname{ScenarioLoaderTool} \\
Trajectory Planner (TP)
  & Trajectory prediction of the flight data system \cite{skybrary_mtcd}
  & Build the nominal 4D trajectory of each flight (position and ETA at each waypoint)
  & \fname{TrajectoryTool} \\
Flight Priority Officer (FPO)
  & Supervisor handling emergency and priority flights \cite{icao_doc4444}
    & Turn the status codes into a priority level and handling rules (for example, do not
    move this aircraft), so that Tactical moves the other aircraft
  & \fname{PriorityRulesTool} \\
\midrule
\multicolumn{4}{@{}l}{\textit{Airspace Constraints crew}} \\
Weather Data Collector (WDC)
  & Meteorological office that supplies data to ATC
  & Load the weather snapshot and check that it covers the whole area, levels and time
  & \fname{WeatherSnapshotTool} \\
Weather Hazard Analyst (WHA)
  & Aviation forecaster issuing hazard warnings (SIGMET) \cite{aviationweather_api}
  & Apply the thresholds to the snapshot to get wind per flight level and hazard zones
    (storms, turbulence, icing), and describe them for the report
  & \fname{WeatherHazardTool} \\
Airspace Manager (AM)
  & Airspace Management Cell \cite{eurocontrol_ernip3}
  & Load restricted and military zones and airport constraints as constraint zones
  & \fname{AirspaceRestrictionTool} \\
\midrule
\multicolumn{4}{@{}l}{\textit{Surveillance \& Conflict crew}} \\
Radar Controller (RC)
  & Controller using conflict detection tools \cite{skybrary_mtcd}
  & Run the detector on all trajectories: aircraft against aircraft and against zones
  & \fname{ConflictDetectionTool} \\
Conflict Analyst (CA)
  & Planner controller \cite{skybrary_exec_planner}
  & Classify each conflict (type, severity, time left, priority flights) and order them
  & \fname{SeverityTool} \\
\midrule
\multicolumn{4}{@{}l}{\textit{Tactical Resolution crew}} \\
Executive Controller (EC)
  & Executive controller \cite{skybrary_exec_planner}
  & Choose one manoeuvre per conflict (level change, speed change, delay or hold,
    reroute), using the feedback of SUP on a retry
  & \fname{ManoeuvreSimulatorTool} \\
\midrule
\multicolumn{4}{@{}l}{\textit{Safety Validation crew}} \\
Verification Officer (VO)
  & Clearance verification \cite{skybrary_mtcd}
  & Run the detector again on the modified flight plans and the constraint zones
  & \fname{ConflictDetectionTool} \\
ATC Supervisor (SUP)
  & ATC shift supervisor
  & Approve or reject each manoeuvre and write structured feedback for EC
  & \fname{SeverityTool} \\
\midrule
\multicolumn{4}{@{}l}{\textit{Reporting crew}} \\
ATC Report Officer (RO)
  & Operational reporting of an ATC unit
  &  Write the Markdown report: situation, conflicts, resolutions, clearance status, flagged
    cases and rejected flight plans
  & \fname{FileReadTool} \\
Radio Communications Officer (RCO)
  & Radiotelephony of the executive controller \cite{skybrary_phraseology, icao_doc9432}
  & Turn each approved manoeuvre into a standard radio instruction
  & \fname{PhraseologyTool} \\
\bottomrule
\end{tabularx}
\end{table}

All the tools are deterministic. \fname{FileReadTool} is a built-in CrewAI tool, and the
other ten are custom tools built with \fname{BaseTool}, as in Lab~4 \cite{crewai_docs}.
\fname{ConflictDetectionTool} already exists: it wraps the detector of the course notebook
(Section~\ref{sec:agentification}) and will be extended to check constraint zones and wind.
Three tools deserve a short comment:

\begin{itemize}
  \item \fname{ManoeuvreSimulatorTool} is a ``what-if'' tool, like the one that supports
        MTCD \cite{skybrary_mtcd}. It applies a candidate manoeuvre to a flight plan,
        rebuilds the trajectory and calls the same detection functions, so EC can test an
        idea before proposing it. VO still checks the final proposal independently.
  \item The feedback of SUP is a Pydantic model with the conflict identifier, the decision,
        the reason (for example, ``new conflict with AC3 at 340~s'' or ``enters storm zone
        Z2'') and a suggestion, so that EC can actually use it on the retry.
  \item \fname{PhraseologyTool} fills fixed templates based on standard phraseology,
        because non-standard phraseology often contributes to incidents
        \cite{skybrary_phraseology}. The agent only chooses the template and the reason, for
        example ``AC2, climb to flight level 320, due traffic''.
\end{itemize}

EC can choose between four types of manoeuvre (Table~\ref{tab:manoeuvres}). Each one changes
the flight plan in a way that \fname{ManoeuvreSimulatorTool} can apply and check, and their
values use fixed steps, so the number of possible actions is finite.

\begin{table}[!htbp]
\centering
\small
\caption{Manoeuvres available to the Executive Controller.}
\label{tab:manoeuvres}
\begin{tabularx}{\textwidth}{@{}>{\raggedright\arraybackslash}p{2.4cm}
                               >{\raggedright\arraybackslash}p{4.6cm} X@{}}
\toprule
\textbf{Manoeuvre} & \textbf{What changes in the flight plan} & \textbf{Values and use} \\
\midrule
Level change
  & Cruise altitude
  & Climb or descend in steps of 1000~ft, within cruise level $\pm$4000~ft (the levels
    covered by the weather snapshot). Usually the first choice for traffic conflicts. \\
\addlinespace
Speed change
  & Speed of the aircraft
  & Steps of 0.01~Mach (about 6~kt at our cruise levels) \cite{skybrary_speed_control},
    at most $\pm$5\% of the planned speed, because at cruise levels the margin between the
    minimum and maximum speed is small \cite{skybrary_cruise_guide}. Not used for weather
    conflicts, because speed control is not reliable during weather avoidance
    \cite{skybrary_wx_los}. \\
\addlinespace
Delay or hold
  & Departure time (ground delay) or extra time at a waypoint (hold)
  & Whole minutes, up to the maximum delay covered by the snapshot. Used for airport
    constraints and when no level or route is free. \\
\addlinespace
Reroute
  & Route
  & Another path of the airspace graph between the same first and last node, so it only
    uses airways that exist. Used to avoid constraint zones. \\
\bottomrule
\end{tabularx}
\end{table}

% -------------------------------------------------
\subsection{System Overview Diagram}
\label{sec:overview}

Figure~\ref{fig:flow} shows the complete Flow: the six crews, the two routers and the
intermediate files.


The two routers are deterministic functions of the Flow \cite{crewai_flows}:

\begin{itemize}
  \item \textbf{Router~1} goes to Tactical Resolution if \fname{02_surveillance.json}
        contains at least one conflict, either between two aircraft or between an aircraft
        and a constraint zone. Otherwise it goes directly to Reporting, as the assignment
        requires \cite{mas_assignment}.
    \item \textbf{Router~2} goes to Reporting if every manoeuvre is approved. The iteration
        counter starts at 0 and counts how many times the flow has gone back to Tactical
        Resolution. If some manoeuvres are rejected and the counter is lower than two, the
        flow goes back to Tactical Resolution with the feedback of SUP, and only the rejected
        conflicts are solved again. If they are still rejected when the counter is two (the
        third attempt), the flow goes to Reporting and the conflict is flagged as ``requires
        human controller'', together with the last manoeuvre and the reason why it was
        rejected.
\end{itemize}

The iteration counter is stored in the state of the Flow and not in any agent, so no
language model can decide to skip the limit.

As an example, in the scenario \fname{sample_conflict} of the course, AC1 and AC2 fly in
opposite directions on the same route at 30000~ft. RC finds a head-on conflict at
$t = 112.5$~s and Router~1 sends it to Tactical. EC proposes that AC2 climbs to FL320. VO
runs the detector again and finds no conflict, because the vertical distance is now
2000~ft, so SUP approves it and Router~2 goes to Reporting. There, RCO writes the
instruction ``AC2, climb to flight level 320, due traffic'' and RO includes it in the
report with the clearance granted.

\begin{figure}[H]
\centering
\begin{tikzpicture}[
    font=\small,
    node distance=6.5mm and 10mm,
    crew/.style={draw, rounded corners=2pt, fill=gray!12, align=center,
                 minimum width=4.8cm, minimum height=9mm},
    extra/.style={crew, fill=blue!12},
    infile/.style={draw, align=center, font=\scriptsize, inner sep=3pt},
    router/.style={draw, diamond, aspect=2.4, fill=yellow!25, align=center,
                   inner sep=0.5pt, font=\scriptsize},
    arr/.style={-{Stealth[length=2mm]}, thick},
    lbl/.style={font=\scriptsize, inner sep=1.5pt}
  ]
  \node[infile, minimum width=4.8cm] (in)
    {\fname{airspace.graphml} \quad \fname{aircraft.json}\\ \fname{flight_plans.json}};
  \node[crew, below=of in] (dm) {Data Management Crew\\{\scriptsize FDA, TP, FPO}};
  \node[extra, below=9mm of dm] (ac) {Airspace Constraints Crew*\\{\scriptsize WDC, WHA, AM}};
  \node[infile, left=10mm of ac] (cin)
    {\fname{weather.json}\\ \fname{restricted_zones.json}\\ \fname{airport_constraints.json}};
  \node[crew, below=9mm of ac] (sc) {Surveillance \& Conflict Crew\\{\scriptsize RC, CA}};
  \node[router, below=of sc] (r1) {Router 1\\conflicts?};
  \node[crew, below=9mm of r1] (tr) {Tactical Resolution Crew\\{\scriptsize EC}};
  \node[extra, below=9mm of tr] (sv) {Safety Validation Crew*\\{\scriptsize VO, SUP}};
  \node[router, below=of sv] (r2) {Router 2\\approved?};
  \node[crew, below=12mm of r2] (rp) {Reporting Crew\\{\scriptsize RO, RCO}};
  \node[infile, below=of rp, double] (out) {\fname{report.md}};

  \draw[arr] (in) -- (dm);
  \draw[arr] (dm) -- node[lbl, right] {\fname{01_data_management.json}} (ac);
  \draw[arr] (cin) -- (ac);
  \draw[arr] (ac) -- node[lbl, right] {\fname{01b_airspace_constraints.json}} (sc);
  \draw[arr] (sc) -- node[lbl, right] {\fname{02_surveillance.json}} (r1);
  \draw[arr] (r1) -- node[lbl, right] {traffic or zone conflicts} (tr);
  \draw[arr] (tr) -- node[lbl, right] {\fname{03_tactical_resolution.json}} (sv);
  \draw[arr] (sv) -- node[lbl, right] {\fname{04_safety_validation.json}} (r2);
  \draw[arr] (r2) -- node[lbl, right, align=left]
    {all approved, or still rejected\\at iteration 2 (flagged)} (rp);
  \draw[arr] (rp) -- (out);

  \draw[arr] (r1.east) -- ++(4.2,0) coordinate (c1)
    |- node[lbl, pos=0.25, right] {no conflicts} (rp.east);
  \draw[arr] (r2.west) -- ++(-2.9,0)
    |- node[lbl, pos=0.25, left, align=right] {rejected,\\iteration $<2$} (tr.west);
\end{tikzpicture}
\caption{CrewAI Flow of the system. Blue crews (*) are the two additional crews. Every crew
can also read the earlier intermediate files; for example, Tactical Resolution and Safety
Validation read the constraint zones of \fname{01b}.}
\label{fig:flow}
\end{figure}

% (it replaces the empty Section 3 headings that are now in main.tex).
% Needs: tikz (arrows.meta, positioning), booktabs, tabularx, array, float -- already in
% the preamble. The diagram also needs shapes.geometric, loaded on the next line.
% =================================================
\usetikzlibrary{shapes.geometric}
\providecommand{\fname}[1]{\nolinkurl{#1}}

\section{Agent Selection and Definition}
\label{sec:agents}

The assignment fixes four crews (Data Management, Surveillance \& Conflict, Tactical
Resolution, and Reporting) and allows each team to add more crews and agents, as long as they are justified \cite{mas_assignment}. We designed the rest of the system together in the course forum and our messaging group.
Each proposal followed a common template (idea, motivation, agents and tools, place in the flow); the rest of the team agreed or answered with concerns, until we reached a solution that satisfied everyone. This section explains the result: which crews we
added and why (Section~\ref{sec:extra-crews}), what each crew does
(Section~\ref{sec:crews}), which agents form each crew (Section~\ref{sec:roles}) and how
everything is connected in the CrewAI Flow (Section~\ref{sec:overview}).
% -------------------------------------------------
\subsection{Decision on Additional Crews}
\label{sec:extra-crews}

Before discussing the proposals, we agreed on five rules to decide about them:

\begin{enumerate}
  \item \textbf{A new crew only for a new step in the flow.} If an idea needs its own input, its own output file and its own decision point, it becomes a crew. If it only adds information to a step that already exists, it becomes an agent inside that crew.
  \item \textbf{One detector.} Every conflict, between two aircraft or between an aircraft and a restricted zone, is found by the same deterministic tool. No crew checks separations with its own logic.
  \item \textbf{Numbers come from tools, decisions come from agents.} Distances, times and thresholds are always computed by deterministic tools. The language model interprets the results, prioritises and decides, but it never computes a separation.
  \item \textbf{Files as the interface.} Each crew reads the files of the previous crews and writes one intermediate file with a Pydantic model, so every crew can be developed and tested on its own, as the assignment recommends \cite{mas_assignment}.
  \item \textbf{Static inputs.} All the information is given as files at the start, and nothing is downloaded while the system runs.
\end{enumerate}

Table~\ref{tab:proposals} summarises the proposals we discussed and what we
decided about each one.

\begin{table}[!htbp]
\centering
\small
\caption{Design proposals and decisions.}
\label{tab:proposals}
\begin{tabularx}{\textwidth}{@{}>{\raggedright\arraybackslash}p{4.3cm}
                               >{\raggedright\arraybackslash}p{3.9cm} X@{}}
\toprule
\textbf{Proposal} & \textbf{Decision} & \textbf{Main reason} \\
\midrule
Meteorological hazards crew
  & New crew, extended to all airspace constraints
  & Wind changes the ETAs and storms create areas to avoid \\
Safety Validation crew
  & New crew, replaces the loop back to Surveillance
  & Tactical should not grade its own work \\
Restricted and military zones (Airspace Manager)
  & Agent in the Airspace Constraints crew
  & Not weather, but the same zone model \\
Airport constraints input
  & New input, handled by the Airspace Manager
  & A closure is a zone placed on a node \\
Flight Priority Officer
  & Agent in Data Management
  & It only adds the priority to the normalised data \\
Radio Communications Officer
  & Agent in Reporting
  & It only changes how the manoeuvres are written \\
Live weather download
  & Rejected: stored snapshot per scenario
  & Repeatable results and tests \\
Weather crew in parallel with Data Management
  & Rejected: sequential
  &  It needs the routes, levels and times from Data Management \\
Runway choice according to wind
  & Future work
  & Approach problem; our data has no runways \\
\bottomrule
\end{tabularx}
\end{table}

\textbf{Airspace Constraints crew.} This crew started as a proposal for meteorological
hazards. Wind changes the ground speed of each aircraft, so it changes the ETAs and
therefore when, and even whether, two aircraft meet. Thunderstorms also create areas that
aircraft must avoid: flight crews try to stay about 20~NM away from a cumulonimbus
\cite{skybrary_wx_avoidance}, and during weather avoidance speed control is not reliable for
separation \cite{skybrary_wx_los}. Later the team proposed two more kinds of restrictions:
restricted or military zones with active hours, which in Europe are managed with the
Flexible Use of Airspace concept \cite{eurocontrol_ernip3}, and airport constraints
(closures, ground stops, reduced capacity). At first we thought of adding the restricted
zones to the weather crew, but they have nothing to do with the weather. Instead of one
crew for each kind of restriction, we created one crew for all the constraints on the
airspace, with a specialised agent for each source. All of them produce the same object, a
\emph{constraint zone}, so the detector and the Safety Validation crew treat every
restriction in the same way.

This crew runs after Data Management, and not in parallel with it, because it needs its
output: the trajectories define the area, the flight levels and the time window that the
constraints must cover, and the first and last node of each route are the airports. It also
has to finish before Surveillance, because the detector needs the wind to compute correct ETAs
and the constraint zones to check the trajectories. The wind is used in two steps:
first the hazard zones, and then the wind-corrected ETAs, which need more changes in the
detector. Both steps are part of the design, but the second one will be implemented only
after the first one works.

\textbf{Safety Validation crew.} A manoeuvre can solve one conflict but create a new one,
or send an aircraft into a hazard zone. In real ATC a clearance is checked before it is
issued, and European controllers have short-term tools for conflict detection and
clearance verification \cite{skybrary_mtcd}. Our first draft of the flow sent the new
flight plans back to the Surveillance crew. We replaced this loop with a separate crew, so
that Tactical does not grade its own work and the feedback has a clear owner. To respect
rule~2, the Verification Officer uses the same tool as Surveillance and reads the
constraint zones from the file of the Airspace Constraints crew, without running that crew
again.

\textbf{Agents instead of crews.} Following rule~1, the other proposals did not need a new
step in the flow, so they became agents in existing crews: the Flight Priority Officer in
Data Management, the Airspace Manager in the Airspace Constraints crew and the Radio
Communications Officer in Reporting.

\textbf{Rejected ideas and future work.} We rejected downloading live weather data during
the execution: every run would give different results, it would be hard to test, and the
assignment gives the inputs as files. Instead, a small script downloads real data from
Open-Meteo \cite{openmeteo_docs} once and saves it in the scenario folder. Choosing the
runway according to the wind is left as future work, because our data has nodes but no
runways, and runway choice belongs to approach control, while our system works en-route
\cite{skybrary_area_controller}.

% -------------------------------------------------
\subsection{Crew Definitions}
\label{sec:crews}

The final system has six crews: the four mandatory ones and two additional ones (marked
with~*). Table~\ref{tab:crews} gives the purpose of each crew, its agents and the files it
reads and writes. All the intermediate files are written in the \fname{output/} folder.

\begin{table}[!htbp]
\centering
\small
\caption{Crews of the system. Files numbered \fname{01}--\fname{04} are intermediate
files in \fname{output/}.}
\label{tab:crews}
\begin{tabularx}{\textwidth}{@{}>{\raggedright\arraybackslash}p{2.3cm} X
                               >{\raggedright\arraybackslash}p{1.25cm}
                               >{\raggedright\arraybackslash}p{3.3cm}
                               >{\raggedright\arraybackslash}p{3.3cm}@{}}
\toprule
\textbf{Crew} & \textbf{Purpose} & \textbf{Agents} & \textbf{Reads} & \textbf{Writes} \\
\midrule
Data Management
  & Load, validate and normalise the inputs, build the nominal 4D trajectories and mark
    priority flights
  & FDA, TP, FPO
  & \fname{airspace.graphml}, \fname{aircraft.json}, \fname{flight_plans.json}
  & \fname{01_data_management.json} \\
\addlinespace
Airspace Constraints*
  & Turn the weather, the restricted zones and the airport constraints into wind per level
    and constraint zones
  & WDC, WHA, AM
  & \fname{01}, \fname{weather.json}, \fname{restricted_zones.json},
    \fname{airport_constraints.json}
  & \fname{01b_airspace_constraints.json} \\
\addlinespace
Surveillance \& Conflict
  & Detect every loss of separation (aircraft--aircraft and aircraft--zone) and prioritise
    the conflicts
  & RC, CA
  & \fname{01}, \fname{01b}
  & \fname{02_surveillance.json} \\
\addlinespace
Tactical Resolution
  & Propose one manoeuvre for each conflict
  & EC
  & \fname{01}, \fname{01b}, \fname{02}, \fname{04} (on a retry)
  & \fname{03_tactical_resolution.json} \\
\addlinespace
Safety Validation*
  & Check every manoeuvre again and approve or reject it with feedback
  & VO, SUP
  & \fname{01}, \fname{01b}, \fname{03}
  & \fname{04_safety_validation.json} \\
\addlinespace
Reporting
  & Write the ATC report and the radio instructions
  & RO, RCO
  & All previous files
  & \fname{report.md} \\
\bottomrule
\end{tabularx}
\end{table}

\textbf{Sequential processes with parallel tasks.} CrewAI also offers a hierarchical process,
where a manager agent delegates the tasks, as we saw in Lab~6 \cite{crewai_docs}. We do not use
it. Inside each crew the order of the tasks is always the same, and the decisions about what
happens next (which crew runs and how many retries are allowed) are taken by the routers of the
Flow, which are deterministic \cite{crewai_flows}. A manager agent would add one more
language-model decision exactly where we want the system to be predictable.

A sequential process does not mean that every task waits for the previous one, though.
Independent tasks run asynchronously, as we also saw in Lab~6. In the Airspace Constraints
crew, the Airspace Manager works in parallel with the weather agents, and in Data Management
the Trajectory Planner and the Flight Priority Officer both start as soon as the Flight Data
Assistant has validated the data. Their results are merged by deterministic code and not by
another agent, since all of them use the same Pydantic models.

This is also why Tactical Resolution has a single agent: one decision maker, checked by a
different crew. As future work, in larger scenarios the conflicts could be split into
independent groups, the connected components of the graph of conflicting aircraft
\cite{durand1995optimal, stollenwerk2017quantum}, and the Flow could run one Tactical Resolution
crew per group in parallel. If a manoeuvre created a new conflict between two groups, the
Safety Validation loop would merge them and solve them again. This
is similar to how real ATC divides the airspace into sectors with a manageable workload
\cite{skybrary_sectorisation}.

\textbf{Inputs.} Each scenario folder contains the three files required by the assignment and
three optional files that we added (Table~\ref{tab:inputs}). All times are given in seconds
from the start of the scenario. If an optional file is missing, the crew writes an empty list,
so the sample scenarios of the course still work. In cruise, aircraft fly at pressure altitude,
so we read the cruise altitude of each flight plan as a flight level (30000~ft = FL300), and
each flight level corresponds to a fixed pressure level in the weather data.

\begin{table}[!htbp]
\centering
\small
\caption{Input files of each scenario.}
\label{tab:inputs}
\begin{tabularx}{\textwidth}{@{}>{\raggedright\arraybackslash}p{3.6cm} X@{}}
\toprule
\textbf{File} & \textbf{Content} \\
\midrule
\multicolumn{2}{@{}l}{\textit{Required by the assignment}} \\
\fname{airspace.graphml}
  & Airspace graph. Each node is a waypoint with its latitude and longitude, and each edge
    is an airway between two waypoints with its distance in NM. \\
\addlinespace
\fname{aircraft.json}
   & One entry per aircraft: identifier and cruise speed in knots. We add an optional
    \fname{status} field with the special-handling indicator of the flight plan (for
    example \fname{STS/HOSP} or \fname{STS/MEDEVAC}) or the emergency transponder code
    \fname{7700} \cite{icao_doc4444}. \\
\addlinespace
\fname{flight_plans.json}
  & One entry per flight: aircraft identifier, route (ordered list of waypoints of the
    graph), departure time and cruise altitude in feet. The first and last waypoints of the
    route are its airports. \\
\midrule
\multicolumn{2}{@{}l}{\textit{Optional, added by us}} \\
\fname{weather.json}
  & Weather snapshot downloaded once from Open-Meteo \cite{openmeteo_docs}: the real date and
    time of the scenario start (\fname{scenario_start_utc}), and a grid of points that covers
    the airspace graph plus a 20~NM margin. For each point and hour, it gives the wind speed
    and direction at each pressure level near the cruise levels, the geopotential height of
    each level (its real altitude that day), the convective available potential energy
    (CAPE) and the freezing level. \\
\addlinespace
\fname{restricted_zones.json}
  & One entry per zone: identifier, type (restricted, military or danger area), shape (a
    circle with centre and radius in NM, or a polygon of latitude/longitude points), lower
    and upper flight level, and the time windows when it is active
    \cite{eurocontrol_ernip3}. \\
\addlinespace
\fname{airport_constraints.json}
  & One entry per constraint: airport node, type (closure, ground stop, reduced capacity,
    works or incident), affected movements (departures, arrivals or both), time window and,
    for reduced capacity, the maximum number of movements per hour. \\
\bottomrule
\end{tabularx}
\end{table}

The thresholds that turn the weather data into hazard zones are kept in one configuration
file, \fname{hazard_thresholds.yaml}, each with its reference: the CAPE above which an area
counts as a convective storm, the buffer around a storm (20~NM \cite{skybrary_wx_avoidance}),
the wind speed above which a level counts as turbulent, and the temperature range at which
icing is likely.

\textbf{Constraint zones.} All the restrictions share one model: an area (a circle or a
polygon, or a graph node for airports), an altitude band, a time window, a type and a
severity. Since the assignment does not define airports, we take the first and the last
node of each route as its airports. Airport constraints have no altitude band: a closure or
a ground stop affects the flights that depart from or arrive at that node inside the time
window, and Tactical solves it with a ground delay or a hold. Only the departure from the first node and the arrival at the last node of a flight plan
count as movements; flights that pass over the node at cruise level are not affected.

\textbf{Invalid flight plans.} If a flight plan refers to an unknown aircraft, a node that is
not in the graph or an airway that does not exist, FDA rejects it and gives the reason, as
the European flight plan processing system does \cite{eurocontrol_ifpl_spec}. The plan is not
repaired, because guessing a route would mean inventing data. The run continues with the
valid flights, and the rejected ones are excluded from detection and resolution, since they
have no trajectory. The report lists them with their reasons and warns that their conflicts
were not checked, because a rejected flight that still has to operate must have its plan
corrected and resubmitted \cite{eurocontrol_ifpl_spec}. If no flight plan is valid, the report
says that there was nothing to check instead of reporting that there are no conflicts.

\textbf{Weather coverage.} The weather snapshot must cover more than the planned routes,
because Tactical can reroute an aircraft, change its level or delay it. For this reason it
covers the whole airspace graph plus a 20~NM margin, every level that Tactical may assign
(cruise level $\pm$4000~ft) and the time window of the scenario including the maximum
allowed delay. In this way the Safety Validation crew never needs new weather data.

% -------------------------------------------------
\subsection{Agent Roles and Tools}
\label{sec:roles}

Each agent is inspired by a real role in air traffic management. As the assignment asks,
we keep the definitions short (role, goal and task) and avoid long backstories.
Table~\ref{tab:roles} defines the abbreviations of the agents used in the rest of the report.

\begin{table}[!htbp]
\centering
\small
\caption{Agents, the real role that inspires them, their main task and their tools.}
\label{tab:roles}
\begin{tabularx}{\textwidth}{@{}>{\raggedright\arraybackslash}p{2.6cm}
                               >{\raggedright\arraybackslash}p{2.75cm} X
                               >{\raggedright\arraybackslash}p{4.25cm}@{}}
\toprule
\textbf{Agent} & \textbf{Inspired by} & \textbf{Main task} & \textbf{Tools} \\
\midrule
\multicolumn{4}{@{}l}{\textit{Data Management crew}} \\
Flight Data Assistant (FDA)
  & Flight data assistant of an area control centre
  & Load the input files, validate the flight plans and normalise units and formats
  & \fname{ScenarioLoaderTool} \\
Trajectory Planner (TP)
  & Trajectory prediction of the flight data system \cite{skybrary_mtcd}
  & Build the nominal 4D trajectory of each flight (position and ETA at each waypoint)
  & \fname{TrajectoryTool} \\
Flight Priority Officer (FPO)
  & Supervisor handling emergency and priority flights \cite{icao_doc4444}
    & Turn the status codes into a priority level and handling rules (for example, do not
    move this aircraft), so that Tactical moves the other aircraft
  & \fname{PriorityRulesTool} \\
\midrule
\multicolumn{4}{@{}l}{\textit{Airspace Constraints crew}} \\
Weather Data Collector (WDC)
  & Meteorological office that supplies data to ATC
  & Load the weather snapshot and check that it covers the whole area, levels and time
  & \fname{WeatherSnapshotTool} \\
Weather Hazard Analyst (WHA)
  & Aviation forecaster issuing hazard warnings (SIGMET) \cite{aviationweather_api}
  & Apply the thresholds to the snapshot to get wind per flight level and hazard zones
    (storms, turbulence, icing), and describe them for the report
  & \fname{WeatherHazardTool} \\
Airspace Manager (AM)
  & Airspace Management Cell \cite{eurocontrol_ernip3}
  & Load restricted and military zones and airport constraints as constraint zones
  & \fname{AirspaceRestrictionTool} \\
\midrule
\multicolumn{4}{@{}l}{\textit{Surveillance \& Conflict crew}} \\
Radar Controller (RC)
  & Controller using conflict detection tools \cite{skybrary_mtcd}
  & Run the detector on all trajectories: aircraft against aircraft and against zones
  & \fname{ConflictDetectionTool} \\
Conflict Analyst (CA)
  & Planner controller \cite{skybrary_exec_planner}
  & Classify each conflict (type, severity, time left, priority flights) and order them
  & \fname{SeverityTool} \\
\midrule
\multicolumn{4}{@{}l}{\textit{Tactical Resolution crew}} \\
Executive Controller (EC)
  & Executive controller \cite{skybrary_exec_planner}
  & Choose one manoeuvre per conflict (level change, speed change, delay or hold,
    reroute), using the feedback of SUP on a retry
  & \fname{ManoeuvreSimulatorTool} \\
\midrule
\multicolumn{4}{@{}l}{\textit{Safety Validation crew}} \\
Verification Officer (VO)
  & Clearance verification \cite{skybrary_mtcd}
  & Run the detector again on the modified flight plans and the constraint zones
  & \fname{ConflictDetectionTool} \\
ATC Supervisor (SUP)
  & ATC shift supervisor
  & Approve or reject each manoeuvre and write structured feedback for EC
  & \fname{SeverityTool} \\
\midrule
\multicolumn{4}{@{}l}{\textit{Reporting crew}} \\
ATC Report Officer (RO)
  & Operational reporting of an ATC unit
  &  Write the Markdown report: situation, conflicts, resolutions, clearance status, flagged
    cases and rejected flight plans
  & \fname{FileReadTool} \\
Radio Communications Officer (RCO)
  & Radiotelephony of the executive controller \cite{skybrary_phraseology, icao_doc9432}
  & Turn each approved manoeuvre into a standard radio instruction
  & \fname{PhraseologyTool} \\
\bottomrule
\end{tabularx}
\end{table}

All the tools are deterministic. \fname{FileReadTool} is a built-in CrewAI tool, and the
other ten are custom tools built with \fname{BaseTool}, as in Lab~4 \cite{crewai_docs}.
\fname{ConflictDetectionTool} already exists: it wraps the detector of the course notebook
(Section~\ref{sec:agentification}) and will be extended to check constraint zones and wind.
Three tools deserve a short comment:

\begin{itemize}
  \item \fname{ManoeuvreSimulatorTool} is a ``what-if'' tool, like the one that supports
        MTCD \cite{skybrary_mtcd}. It applies a candidate manoeuvre to a flight plan,
        rebuilds the trajectory and calls the same detection functions, so EC can test an
        idea before proposing it. VO still checks the final proposal independently.
  \item The feedback of SUP is a Pydantic model with the conflict identifier, the decision,
        the reason (for example, ``new conflict with AC3 at 340~s'' or ``enters storm zone
        Z2'') and a suggestion, so that EC can actually use it on the retry.
  \item \fname{PhraseologyTool} fills fixed templates based on standard phraseology,
        because non-standard phraseology often contributes to incidents
        \cite{skybrary_phraseology}. The agent only chooses the template and the reason, for
        example ``AC2, climb to flight level 320, due traffic''.
\end{itemize}

EC can choose between four types of manoeuvre (Table~\ref{tab:manoeuvres}). Each one changes
the flight plan in a way that \fname{ManoeuvreSimulatorTool} can apply and check, and their
values use fixed steps, so the number of possible actions is finite.

\begin{table}[!htbp]
\centering
\small
\caption{Manoeuvres available to the Executive Controller.}
\label{tab:manoeuvres}
\begin{tabularx}{\textwidth}{@{}>{\raggedright\arraybackslash}p{2.4cm}
                               >{\raggedright\arraybackslash}p{4.6cm} X@{}}
\toprule
\textbf{Manoeuvre} & \textbf{What changes in the flight plan} & \textbf{Values and use} \\
\midrule
Level change
  & Cruise altitude
  & Climb or descend in steps of 1000~ft, within cruise level $\pm$4000~ft (the levels
    covered by the weather snapshot). Usually the first choice for traffic conflicts. \\
\addlinespace
Speed change
  & Speed of the aircraft
  & Steps of 0.01~Mach (about 6~kt at our cruise levels) \cite{skybrary_speed_control},
    at most $\pm$5\% of the planned speed, because at cruise levels the margin between the
    minimum and maximum speed is small \cite{skybrary_cruise_guide}. Not used for weather
    conflicts, because speed control is not reliable during weather avoidance
    \cite{skybrary_wx_los}. \\
\addlinespace
Delay or hold
  & Departure time (ground delay) or extra time at a waypoint (hold)
  & Whole minutes, up to the maximum delay covered by the snapshot. Used for airport
    constraints and when no level or route is free. \\
\addlinespace
Reroute
  & Route
  & Another path of the airspace graph between the same first and last node, so it only
    uses airways that exist. Used to avoid constraint zones. \\
\bottomrule
\end{tabularx}
\end{table}

% -------------------------------------------------
\subsection{System Overview Diagram}
\label{sec:overview}

Figure~\ref{fig:flow} shows the complete Flow: the six crews, the two routers and the
intermediate files.


The two routers are deterministic functions of the Flow \cite{crewai_flows}:

\begin{itemize}
  \item \textbf{Router~1} goes to Tactical Resolution if \fname{02_surveillance.json}
        contains at least one conflict, either between two aircraft or between an aircraft
        and a constraint zone. Otherwise it goes directly to Reporting, as the assignment
        requires \cite{mas_assignment}.
    \item \textbf{Router~2} goes to Reporting if every manoeuvre is approved. The iteration
        counter starts at 0 and counts how many times the flow has gone back to Tactical
        Resolution. If some manoeuvres are rejected and the counter is lower than two, the
        flow goes back to Tactical Resolution with the feedback of SUP, and only the rejected
        conflicts are solved again. If they are still rejected when the counter is two (the
        third attempt), the flow goes to Reporting and the conflict is flagged as ``requires
        human controller'', together with the last manoeuvre and the reason why it was
        rejected.
\end{itemize}

The iteration counter is stored in the state of the Flow and not in any agent, so no
language model can decide to skip the limit.

As an example, in the scenario \fname{sample_conflict} of the course, AC1 and AC2 fly in
opposite directions on the same route at 30000~ft. RC finds a head-on conflict at
$t = 112.5$~s and Router~1 sends it to Tactical. EC proposes that AC2 climbs to FL320. VO
runs the detector again and finds no conflict, because the vertical distance is now
2000~ft, so SUP approves it and Router~2 goes to Reporting. There, RCO writes the
instruction ``AC2, climb to flight level 320, due traffic'' and RO includes it in the
report with the clearance granted.

\begin{figure}[H]
\centering
\begin{tikzpicture}[
    font=\small,
    node distance=6.5mm and 10mm,
    crew/.style={draw, rounded corners=2pt, fill=gray!12, align=center,
                 minimum width=4.8cm, minimum height=9mm},
    extra/.style={crew, fill=blue!12},
    infile/.style={draw, align=center, font=\scriptsize, inner sep=3pt},
    router/.style={draw, diamond, aspect=2.4, fill=yellow!25, align=center,
                   inner sep=0.5pt, font=\scriptsize},
    arr/.style={-{Stealth[length=2mm]}, thick},
    lbl/.style={font=\scriptsize, inner sep=1.5pt}
  ]
  \node[infile, minimum width=4.8cm] (in)
    {\fname{airspace.graphml} \quad \fname{aircraft.json}\\ \fname{flight_plans.json}};
  \node[crew, below=of in] (dm) {Data Management Crew\\{\scriptsize FDA, TP, FPO}};
  \node[extra, below=9mm of dm] (ac) {Airspace Constraints Crew*\\{\scriptsize WDC, WHA, AM}};
  \node[infile, left=10mm of ac] (cin)
    {\fname{weather.json}\\ \fname{restricted_zones.json}\\ \fname{airport_constraints.json}};
  \node[crew, below=9mm of ac] (sc) {Surveillance \& Conflict Crew\\{\scriptsize RC, CA}};
  \node[router, below=of sc] (r1) {Router 1\\conflicts?};
  \node[crew, below=9mm of r1] (tr) {Tactical Resolution Crew\\{\scriptsize EC}};
  \node[extra, below=9mm of tr] (sv) {Safety Validation Crew*\\{\scriptsize VO, SUP}};
  \node[router, below=of sv] (r2) {Router 2\\approved?};
  \node[crew, below=12mm of r2] (rp) {Reporting Crew\\{\scriptsize RO, RCO}};
  \node[infile, below=of rp, double] (out) {\fname{report.md}};

  \draw[arr] (in) -- (dm);
  \draw[arr] (dm) -- node[lbl, right] {\fname{01_data_management.json}} (ac);
  \draw[arr] (cin) -- (ac);
  \draw[arr] (ac) -- node[lbl, right] {\fname{01b_airspace_constraints.json}} (sc);
  \draw[arr] (sc) -- node[lbl, right] {\fname{02_surveillance.json}} (r1);
  \draw[arr] (r1) -- node[lbl, right] {traffic or zone conflicts} (tr);
  \draw[arr] (tr) -- node[lbl, right] {\fname{03_tactical_resolution.json}} (sv);
  \draw[arr] (sv) -- node[lbl, right] {\fname{04_safety_validation.json}} (r2);
  \draw[arr] (r2) -- node[lbl, right, align=left]
    {all approved, or still rejected\\at iteration 2 (flagged)} (rp);
  \draw[arr] (rp) -- (out);

  \draw[arr] (r1.east) -- ++(4.2,0) coordinate (c1)
    |- node[lbl, pos=0.25, right] {no conflicts} (rp.east);
  \draw[arr] (r2.west) -- ++(-2.9,0)
    |- node[lbl, pos=0.25, left, align=right] {rejected,\\iteration $<2$} (tr.west);
\end{tikzpicture}
\caption{CrewAI Flow of the system. Blue crews (*) are the two additional crews. Every crew
can also read the earlier intermediate files; for example, Tactical Resolution and Safety
Validation read the constraint zones of \fname{01b}.}
\label{fig:flow}
\end{figure}

% (it replaces the empty Section 3 headings that are now in main.tex).
% Needs: tikz (arrows.meta, positioning), booktabs, tabularx, array, float -- already in
% the preamble. The diagram also needs shapes.geometric, loaded on the next line.
% =================================================
\usetikzlibrary{shapes.geometric}
\providecommand{\fname}[1]{\nolinkurl{#1}}

\section{Agent Selection and Definition}
\label{sec:agents}

The assignment fixes four crews (Data Management, Surveillance \& Conflict, Tactical
Resolution, and Reporting) and allows each team to add more crews and agents, as long as they are justified \cite{mas_assignment}. We designed the rest of the system together in the course forum and our messaging group.
Each proposal followed a common template (idea, motivation, agents and tools, place in the flow); the rest of the team agreed or answered with concerns, until we reached a solution that satisfied everyone. This section explains the result: which crews we
added and why (Section~\ref{sec:extra-crews}), what each crew does
(Section~\ref{sec:crews}), which agents form each crew (Section~\ref{sec:roles}) and how
everything is connected in the CrewAI Flow (Section~\ref{sec:overview}).
% -------------------------------------------------
\subsection{Decision on Additional Crews}
\label{sec:extra-crews}

Before discussing the proposals, we agreed on five rules to decide about them:

\begin{enumerate}
  \item \textbf{A new crew only for a new step in the flow.} If an idea needs its own input, its own output file and its own decision point, it becomes a crew. If it only adds information to a step that already exists, it becomes an agent inside that crew.
  \item \textbf{One detector.} Every conflict, between two aircraft or between an aircraft and a restricted zone, is found by the same deterministic tool. No crew checks separations with its own logic.
  \item \textbf{Numbers come from tools, decisions come from agents.} Distances, times and thresholds are always computed by deterministic tools. The language model interprets the results, prioritises and decides, but it never computes a separation.
  \item \textbf{Files as the interface.} Each crew reads the files of the previous crews and writes one intermediate file with a Pydantic model, so every crew can be developed and tested on its own, as the assignment recommends \cite{mas_assignment}.
  \item \textbf{Static inputs.} All the information is given as files at the start, and nothing is downloaded while the system runs.
\end{enumerate}

Table~\ref{tab:proposals} summarises the proposals we discussed and what we
decided about each one.

\begin{table}[!htbp]
\centering
\small
\caption{Design proposals and decisions.}
\label{tab:proposals}
\begin{tabularx}{\textwidth}{@{}>{\raggedright\arraybackslash}p{4.3cm}
                               >{\raggedright\arraybackslash}p{3.9cm} X@{}}
\toprule
\textbf{Proposal} & \textbf{Decision} & \textbf{Main reason} \\
\midrule
Meteorological hazards crew
  & New crew, extended to all airspace constraints
  & Wind changes the ETAs and storms create areas to avoid \\
Safety Validation crew
  & New crew, replaces the loop back to Surveillance
  & Tactical should not grade its own work \\
Restricted and military zones (Airspace Manager)
  & Agent in the Airspace Constraints crew
  & Not weather, but the same zone model \\
Airport constraints input
  & New input, handled by the Airspace Manager
  & A closure is a zone placed on a node \\
Flight Priority Officer
  & Agent in Data Management
  & It only adds the priority to the normalised data \\
Radio Communications Officer
  & Agent in Reporting
  & It only changes how the manoeuvres are written \\
Live weather download
  & Rejected: stored snapshot per scenario
  & Repeatable results and tests \\
Weather crew in parallel with Data Management
  & Rejected: sequential
  &  It needs the routes, levels and times from Data Management \\
Runway choice according to wind
  & Future work
  & Approach problem; our data has no runways \\
\bottomrule
\end{tabularx}
\end{table}

\textbf{Airspace Constraints crew.} This crew started as a proposal for meteorological
hazards. Wind changes the ground speed of each aircraft, so it changes the ETAs and
therefore when, and even whether, two aircraft meet. Thunderstorms also create areas that
aircraft must avoid: flight crews try to stay about 20~NM away from a cumulonimbus
\cite{skybrary_wx_avoidance}, and during weather avoidance speed control is not reliable for
separation \cite{skybrary_wx_los}. Later the team proposed two more kinds of restrictions:
restricted or military zones with active hours, which in Europe are managed with the
Flexible Use of Airspace concept \cite{eurocontrol_ernip3}, and airport constraints
(closures, ground stops, reduced capacity). At first we thought of adding the restricted
zones to the weather crew, but they have nothing to do with the weather. Instead of one
crew for each kind of restriction, we created one crew for all the constraints on the
airspace, with a specialised agent for each source. All of them produce the same object, a
\emph{constraint zone}, so the detector and the Safety Validation crew treat every
restriction in the same way.

This crew runs after Data Management, and not in parallel with it, because it needs its
output: the trajectories define the area, the flight levels and the time window that the
constraints must cover, and the first and last node of each route are the airports. It also
has to finish before Surveillance, because the detector needs the wind to compute correct ETAs
and the constraint zones to check the trajectories. The wind is used in two steps:
first the hazard zones, and then the wind-corrected ETAs, which need more changes in the
detector. Both steps are part of the design, but the second one will be implemented only
after the first one works.

\textbf{Safety Validation crew.} A manoeuvre can solve one conflict but create a new one,
or send an aircraft into a hazard zone. In real ATC a clearance is checked before it is
issued, and European controllers have short-term tools for conflict detection and
clearance verification \cite{skybrary_mtcd}. Our first draft of the flow sent the new
flight plans back to the Surveillance crew. We replaced this loop with a separate crew, so
that Tactical does not grade its own work and the feedback has a clear owner. To respect
rule~2, the Verification Officer uses the same tool as Surveillance and reads the
constraint zones from the file of the Airspace Constraints crew, without running that crew
again.

\textbf{Agents instead of crews.} Following rule~1, the other proposals did not need a new
step in the flow, so they became agents in existing crews: the Flight Priority Officer in
Data Management, the Airspace Manager in the Airspace Constraints crew and the Radio
Communications Officer in Reporting.

\textbf{Rejected ideas and future work.} We rejected downloading live weather data during
the execution: every run would give different results, it would be hard to test, and the
assignment gives the inputs as files. Instead, a small script downloads real data from
Open-Meteo \cite{openmeteo_docs} once and saves it in the scenario folder. Choosing the
runway according to the wind is left as future work, because our data has nodes but no
runways, and runway choice belongs to approach control, while our system works en-route
\cite{skybrary_area_controller}.

% -------------------------------------------------
\subsection{Crew Definitions}
\label{sec:crews}

The final system has six crews: the four mandatory ones and two additional ones (marked
with~*). Table~\ref{tab:crews} gives the purpose of each crew, its agents and the files it
reads and writes. All the intermediate files are written in the \fname{output/} folder.

\begin{table}[!htbp]
\centering
\small
\caption{Crews of the system. Files numbered \fname{01}--\fname{04} are intermediate
files in \fname{output/}.}
\label{tab:crews}
\begin{tabularx}{\textwidth}{@{}>{\raggedright\arraybackslash}p{2.3cm} X
                               >{\raggedright\arraybackslash}p{1.25cm}
                               >{\raggedright\arraybackslash}p{3.3cm}
                               >{\raggedright\arraybackslash}p{3.3cm}@{}}
\toprule
\textbf{Crew} & \textbf{Purpose} & \textbf{Agents} & \textbf{Reads} & \textbf{Writes} \\
\midrule
Data Management
  & Load, validate and normalise the inputs, build the nominal 4D trajectories and mark
    priority flights
  & FDA, TP, FPO
  & \fname{airspace.graphml}, \fname{aircraft.json}, \fname{flight_plans.json}
  & \fname{01_data_management.json} \\
\addlinespace
Airspace Constraints*
  & Turn the weather, the restricted zones and the airport constraints into wind per level
    and constraint zones
  & WDC, WHA, AM
  & \fname{01}, \fname{weather.json}, \fname{restricted_zones.json},
    \fname{airport_constraints.json}
  & \fname{01b_airspace_constraints.json} \\
\addlinespace
Surveillance \& Conflict
  & Detect every loss of separation (aircraft--aircraft and aircraft--zone) and prioritise
    the conflicts
  & RC, CA
  & \fname{01}, \fname{01b}
  & \fname{02_surveillance.json} \\
\addlinespace
Tactical Resolution
  & Propose one manoeuvre for each conflict
  & EC
  & \fname{01}, \fname{01b}, \fname{02}, \fname{04} (on a retry)
  & \fname{03_tactical_resolution.json} \\
\addlinespace
Safety Validation*
  & Check every manoeuvre again and approve or reject it with feedback
  & VO, SUP
  & \fname{01}, \fname{01b}, \fname{03}
  & \fname{04_safety_validation.json} \\
\addlinespace
Reporting
  & Write the ATC report and the radio instructions
  & RO, RCO
  & All previous files
  & \fname{report.md} \\
\bottomrule
\end{tabularx}
\end{table}

\textbf{Sequential processes with parallel tasks.} CrewAI also offers a hierarchical process,
where a manager agent delegates the tasks, as we saw in Lab~6 \cite{crewai_docs}. We do not use
it. Inside each crew the order of the tasks is always the same, and the decisions about what
happens next (which crew runs and how many retries are allowed) are taken by the routers of the
Flow, which are deterministic \cite{crewai_flows}. A manager agent would add one more
language-model decision exactly where we want the system to be predictable.

A sequential process does not mean that every task waits for the previous one, though.
Independent tasks run asynchronously, as we also saw in Lab~6. In the Airspace Constraints
crew, the Airspace Manager works in parallel with the weather agents, and in Data Management
the Trajectory Planner and the Flight Priority Officer both start as soon as the Flight Data
Assistant has validated the data. Their results are merged by deterministic code and not by
another agent, since all of them use the same Pydantic models.

This is also why Tactical Resolution has a single agent: one decision maker, checked by a
different crew. As future work, in larger scenarios the conflicts could be split into
independent groups, the connected components of the graph of conflicting aircraft
\cite{durand1995optimal, stollenwerk2017quantum}, and the Flow could run one Tactical Resolution
crew per group in parallel. If a manoeuvre created a new conflict between two groups, the
Safety Validation loop would merge them and solve them again. This
is similar to how real ATC divides the airspace into sectors with a manageable workload
\cite{skybrary_sectorisation}.

\textbf{Inputs.} Each scenario folder contains the three files required by the assignment and
three optional files that we added (Table~\ref{tab:inputs}). All times are given in seconds
from the start of the scenario. If an optional file is missing, the crew writes an empty list,
so the sample scenarios of the course still work. In cruise, aircraft fly at pressure altitude,
so we read the cruise altitude of each flight plan as a flight level (30000~ft = FL300), and
each flight level corresponds to a fixed pressure level in the weather data.

\begin{table}[!htbp]
\centering
\small
\caption{Input files of each scenario.}
\label{tab:inputs}
\begin{tabularx}{\textwidth}{@{}>{\raggedright\arraybackslash}p{3.6cm} X@{}}
\toprule
\textbf{File} & \textbf{Content} \\
\midrule
\multicolumn{2}{@{}l}{\textit{Required by the assignment}} \\
\fname{airspace.graphml}
  & Airspace graph. Each node is a waypoint with its latitude and longitude, and each edge
    is an airway between two waypoints with its distance in NM. \\
\addlinespace
\fname{aircraft.json}
   & One entry per aircraft: identifier and cruise speed in knots. We add an optional
    \fname{status} field with the special-handling indicator of the flight plan (for
    example \fname{STS/HOSP} or \fname{STS/MEDEVAC}) or the emergency transponder code
    \fname{7700} \cite{icao_doc4444}. \\
\addlinespace
\fname{flight_plans.json}
  & One entry per flight: aircraft identifier, route (ordered list of waypoints of the
    graph), departure time and cruise altitude in feet. The first and last waypoints of the
    route are its airports. \\
\midrule
\multicolumn{2}{@{}l}{\textit{Optional, added by us}} \\
\fname{weather.json}
  & Weather snapshot downloaded once from Open-Meteo \cite{openmeteo_docs}: the real date and
    time of the scenario start (\fname{scenario_start_utc}), and a grid of points that covers
    the airspace graph plus a 20~NM margin. For each point and hour, it gives the wind speed
    and direction at each pressure level near the cruise levels, the geopotential height of
    each level (its real altitude that day), the convective available potential energy
    (CAPE) and the freezing level. \\
\addlinespace
\fname{restricted_zones.json}
  & One entry per zone: identifier, type (restricted, military or danger area), shape (a
    circle with centre and radius in NM, or a polygon of latitude/longitude points), lower
    and upper flight level, and the time windows when it is active
    \cite{eurocontrol_ernip3}. \\
\addlinespace
\fname{airport_constraints.json}
  & One entry per constraint: airport node, type (closure, ground stop, reduced capacity,
    works or incident), affected movements (departures, arrivals or both), time window and,
    for reduced capacity, the maximum number of movements per hour. \\
\bottomrule
\end{tabularx}
\end{table}

The thresholds that turn the weather data into hazard zones are kept in one configuration
file, \fname{hazard_thresholds.yaml}, each with its reference: the CAPE above which an area
counts as a convective storm, the buffer around a storm (20~NM \cite{skybrary_wx_avoidance}),
the wind speed above which a level counts as turbulent, and the temperature range at which
icing is likely.

\textbf{Constraint zones.} All the restrictions share one model: an area (a circle or a
polygon, or a graph node for airports), an altitude band, a time window, a type and a
severity. Since the assignment does not define airports, we take the first and the last
node of each route as its airports. Airport constraints have no altitude band: a closure or
a ground stop affects the flights that depart from or arrive at that node inside the time
window, and Tactical solves it with a ground delay or a hold. Only the departure from the first node and the arrival at the last node of a flight plan
count as movements; flights that pass over the node at cruise level are not affected.

\textbf{Invalid flight plans.} If a flight plan refers to an unknown aircraft, a node that is
not in the graph or an airway that does not exist, FDA rejects it and gives the reason, as
the European flight plan processing system does \cite{eurocontrol_ifpl_spec}. The plan is not
repaired, because guessing a route would mean inventing data. The run continues with the
valid flights, and the rejected ones are excluded from detection and resolution, since they
have no trajectory. The report lists them with their reasons and warns that their conflicts
were not checked, because a rejected flight that still has to operate must have its plan
corrected and resubmitted \cite{eurocontrol_ifpl_spec}. If no flight plan is valid, the report
says that there was nothing to check instead of reporting that there are no conflicts.

\textbf{Weather coverage.} The weather snapshot must cover more than the planned routes,
because Tactical can reroute an aircraft, change its level or delay it. For this reason it
covers the whole airspace graph plus a 20~NM margin, every level that Tactical may assign
(cruise level $\pm$4000~ft) and the time window of the scenario including the maximum
allowed delay. In this way the Safety Validation crew never needs new weather data.

% -------------------------------------------------
\subsection{Agent Roles and Tools}
\label{sec:roles}

Each agent is inspired by a real role in air traffic management. As the assignment asks,
we keep the definitions short (role, goal and task) and avoid long backstories.
Table~\ref{tab:roles} defines the abbreviations of the agents used in the rest of the report.

\begin{table}[!htbp]
\centering
\small
\caption{Agents, the real role that inspires them, their main task and their tools.}
\label{tab:roles}
\begin{tabularx}{\textwidth}{@{}>{\raggedright\arraybackslash}p{2.6cm}
                               >{\raggedright\arraybackslash}p{2.75cm} X
                               >{\raggedright\arraybackslash}p{4.25cm}@{}}
\toprule
\textbf{Agent} & \textbf{Inspired by} & \textbf{Main task} & \textbf{Tools} \\
\midrule
\multicolumn{4}{@{}l}{\textit{Data Management crew}} \\
Flight Data Assistant (FDA)
  & Flight data assistant of an area control centre
  & Load the input files, validate the flight plans and normalise units and formats
  & \fname{ScenarioLoaderTool} \\
Trajectory Planner (TP)
  & Trajectory prediction of the flight data system \cite{skybrary_mtcd}
  & Build the nominal 4D trajectory of each flight (position and ETA at each waypoint)
  & \fname{TrajectoryTool} \\
Flight Priority Officer (FPO)
  & Supervisor handling emergency and priority flights \cite{icao_doc4444}
    & Turn the status codes into a priority level and handling rules (for example, do not
    move this aircraft), so that Tactical moves the other aircraft
  & \fname{PriorityRulesTool} \\
\midrule
\multicolumn{4}{@{}l}{\textit{Airspace Constraints crew}} \\
Weather Data Collector (WDC)
  & Meteorological office that supplies data to ATC
  & Load the weather snapshot and check that it covers the whole area, levels and time
  & \fname{WeatherSnapshotTool} \\
Weather Hazard Analyst (WHA)
  & Aviation forecaster issuing hazard warnings (SIGMET) \cite{aviationweather_api}
  & Apply the thresholds to the snapshot to get wind per flight level and hazard zones
    (storms, turbulence, icing), and describe them for the report
  & \fname{WeatherHazardTool} \\
Airspace Manager (AM)
  & Airspace Management Cell \cite{eurocontrol_ernip3}
  & Load restricted and military zones and airport constraints as constraint zones
  & \fname{AirspaceRestrictionTool} \\
\midrule
\multicolumn{4}{@{}l}{\textit{Surveillance \& Conflict crew}} \\
Radar Controller (RC)
  & Controller using conflict detection tools \cite{skybrary_mtcd}
  & Run the detector on all trajectories: aircraft against aircraft and against zones
  & \fname{ConflictDetectionTool} \\
Conflict Analyst (CA)
  & Planner controller \cite{skybrary_exec_planner}
  & Classify each conflict (type, severity, time left, priority flights) and order them
  & \fname{SeverityTool} \\
\midrule
\multicolumn{4}{@{}l}{\textit{Tactical Resolution crew}} \\
Executive Controller (EC)
  & Executive controller \cite{skybrary_exec_planner}
  & Choose one manoeuvre per conflict (level change, speed change, delay or hold,
    reroute), using the feedback of SUP on a retry
  & \fname{ManoeuvreSimulatorTool} \\
\midrule
\multicolumn{4}{@{}l}{\textit{Safety Validation crew}} \\
Verification Officer (VO)
  & Clearance verification \cite{skybrary_mtcd}
  & Run the detector again on the modified flight plans and the constraint zones
  & \fname{ConflictDetectionTool} \\
ATC Supervisor (SUP)
  & ATC shift supervisor
  & Approve or reject each manoeuvre and write structured feedback for EC
  & \fname{SeverityTool} \\
\midrule
\multicolumn{4}{@{}l}{\textit{Reporting crew}} \\
ATC Report Officer (RO)
  & Operational reporting of an ATC unit
  &  Write the Markdown report: situation, conflicts, resolutions, clearance status, flagged
    cases and rejected flight plans
  & \fname{FileReadTool} \\
Radio Communications Officer (RCO)
  & Radiotelephony of the executive controller \cite{skybrary_phraseology, icao_doc9432}
  & Turn each approved manoeuvre into a standard radio instruction
  & \fname{PhraseologyTool} \\
\bottomrule
\end{tabularx}
\end{table}

All the tools are deterministic. \fname{FileReadTool} is a built-in CrewAI tool, and the
other ten are custom tools built with \fname{BaseTool}, as in Lab~4 \cite{crewai_docs}.
\fname{ConflictDetectionTool} already exists: it wraps the detector of the course notebook
(Section~\ref{sec:agentification}) and will be extended to check constraint zones and wind.
Three tools deserve a short comment:

\begin{itemize}
  \item \fname{ManoeuvreSimulatorTool} is a ``what-if'' tool, like the one that supports
        MTCD \cite{skybrary_mtcd}. It applies a candidate manoeuvre to a flight plan,
        rebuilds the trajectory and calls the same detection functions, so EC can test an
        idea before proposing it. VO still checks the final proposal independently.
  \item The feedback of SUP is a Pydantic model with the conflict identifier, the decision,
        the reason (for example, ``new conflict with AC3 at 340~s'' or ``enters storm zone
        Z2'') and a suggestion, so that EC can actually use it on the retry.
  \item \fname{PhraseologyTool} fills fixed templates based on standard phraseology,
        because non-standard phraseology often contributes to incidents
        \cite{skybrary_phraseology}. The agent only chooses the template and the reason, for
        example ``AC2, climb to flight level 320, due traffic''.
\end{itemize}

EC can choose between four types of manoeuvre (Table~\ref{tab:manoeuvres}). Each one changes
the flight plan in a way that \fname{ManoeuvreSimulatorTool} can apply and check, and their
values use fixed steps, so the number of possible actions is finite.

\begin{table}[!htbp]
\centering
\small
\caption{Manoeuvres available to the Executive Controller.}
\label{tab:manoeuvres}
\begin{tabularx}{\textwidth}{@{}>{\raggedright\arraybackslash}p{2.4cm}
                               >{\raggedright\arraybackslash}p{4.6cm} X@{}}
\toprule
\textbf{Manoeuvre} & \textbf{What changes in the flight plan} & \textbf{Values and use} \\
\midrule
Level change
  & Cruise altitude
  & Climb or descend in steps of 1000~ft, within cruise level $\pm$4000~ft (the levels
    covered by the weather snapshot). Usually the first choice for traffic conflicts. \\
\addlinespace
Speed change
  & Speed of the aircraft
  & Steps of 0.01~Mach (about 6~kt at our cruise levels) \cite{skybrary_speed_control},
    at most $\pm$5\% of the planned speed, because at cruise levels the margin between the
    minimum and maximum speed is small \cite{skybrary_cruise_guide}. Not used for weather
    conflicts, because speed control is not reliable during weather avoidance
    \cite{skybrary_wx_los}. \\
\addlinespace
Delay or hold
  & Departure time (ground delay) or extra time at a waypoint (hold)
  & Whole minutes, up to the maximum delay covered by the snapshot. Used for airport
    constraints and when no level or route is free. \\
\addlinespace
Reroute
  & Route
  & Another path of the airspace graph between the same first and last node, so it only
    uses airways that exist. Used to avoid constraint zones. \\
\bottomrule
\end{tabularx}
\end{table}

% -------------------------------------------------
\subsection{System Overview Diagram}
\label{sec:overview}

Figure~\ref{fig:flow} shows the complete Flow: the six crews, the two routers and the
intermediate files.


The two routers are deterministic functions of the Flow \cite{crewai_flows}:

\begin{itemize}
  \item \textbf{Router~1} goes to Tactical Resolution if \fname{02_surveillance.json}
        contains at least one conflict, either between two aircraft or between an aircraft
        and a constraint zone. Otherwise it goes directly to Reporting, as the assignment
        requires \cite{mas_assignment}.
    \item \textbf{Router~2} goes to Reporting if every manoeuvre is approved. The iteration
        counter starts at 0 and counts how many times the flow has gone back to Tactical
        Resolution. If some manoeuvres are rejected and the counter is lower than two, the
        flow goes back to Tactical Resolution with the feedback of SUP, and only the rejected
        conflicts are solved again. If they are still rejected when the counter is two (the
        third attempt), the flow goes to Reporting and the conflict is flagged as ``requires
        human controller'', together with the last manoeuvre and the reason why it was
        rejected.
\end{itemize}

The iteration counter is stored in the state of the Flow and not in any agent, so no
language model can decide to skip the limit.

As an example, in the scenario \fname{sample_conflict} of the course, AC1 and AC2 fly in
opposite directions on the same route at 30000~ft. RC finds a head-on conflict at
$t = 112.5$~s and Router~1 sends it to Tactical. EC proposes that AC2 climbs to FL320. VO
runs the detector again and finds no conflict, because the vertical distance is now
2000~ft, so SUP approves it and Router~2 goes to Reporting. There, RCO writes the
instruction ``AC2, climb to flight level 320, due traffic'' and RO includes it in the
report with the clearance granted.

\begin{figure}[H]
\centering
\begin{tikzpicture}[
    font=\small,
    node distance=6.5mm and 10mm,
    crew/.style={draw, rounded corners=2pt, fill=gray!12, align=center,
                 minimum width=4.8cm, minimum height=9mm},
    extra/.style={crew, fill=blue!12},
    infile/.style={draw, align=center, font=\scriptsize, inner sep=3pt},
    router/.style={draw, diamond, aspect=2.4, fill=yellow!25, align=center,
                   inner sep=0.5pt, font=\scriptsize},
    arr/.style={-{Stealth[length=2mm]}, thick},
    lbl/.style={font=\scriptsize, inner sep=1.5pt}
  ]
  \node[infile, minimum width=4.8cm] (in)
    {\fname{airspace.graphml} \quad \fname{aircraft.json}\\ \fname{flight_plans.json}};
  \node[crew, below=of in] (dm) {Data Management Crew\\{\scriptsize FDA, TP, FPO}};
  \node[extra, below=9mm of dm] (ac) {Airspace Constraints Crew*\\{\scriptsize WDC, WHA, AM}};
  \node[infile, left=10mm of ac] (cin)
    {\fname{weather.json}\\ \fname{restricted_zones.json}\\ \fname{airport_constraints.json}};
  \node[crew, below=9mm of ac] (sc) {Surveillance \& Conflict Crew\\{\scriptsize RC, CA}};
  \node[router, below=of sc] (r1) {Router 1\\conflicts?};
  \node[crew, below=9mm of r1] (tr) {Tactical Resolution Crew\\{\scriptsize EC}};
  \node[extra, below=9mm of tr] (sv) {Safety Validation Crew*\\{\scriptsize VO, SUP}};
  \node[router, below=of sv] (r2) {Router 2\\approved?};
  \node[crew, below=12mm of r2] (rp) {Reporting Crew\\{\scriptsize RO, RCO}};
  \node[infile, below=of rp, double] (out) {\fname{report.md}};

  \draw[arr] (in) -- (dm);
  \draw[arr] (dm) -- node[lbl, right] {\fname{01_data_management.json}} (ac);
  \draw[arr] (cin) -- (ac);
  \draw[arr] (ac) -- node[lbl, right] {\fname{01b_airspace_constraints.json}} (sc);
  \draw[arr] (sc) -- node[lbl, right] {\fname{02_surveillance.json}} (r1);
  \draw[arr] (r1) -- node[lbl, right] {traffic or zone conflicts} (tr);
  \draw[arr] (tr) -- node[lbl, right] {\fname{03_tactical_resolution.json}} (sv);
  \draw[arr] (sv) -- node[lbl, right] {\fname{04_safety_validation.json}} (r2);
  \draw[arr] (r2) -- node[lbl, right, align=left]
    {all approved, or still rejected\\at iteration 2 (flagged)} (rp);
  \draw[arr] (rp) -- (out);

  \draw[arr] (r1.east) -- ++(4.2,0) coordinate (c1)
    |- node[lbl, pos=0.25, right] {no conflicts} (rp.east);
  \draw[arr] (r2.west) -- ++(-2.9,0)
    |- node[lbl, pos=0.25, left, align=right] {rejected,\\iteration $<2$} (tr.west);
\end{tikzpicture}
\caption{CrewAI Flow of the system. Blue crews (*) are the two additional crews. Every crew
can also read the earlier intermediate files; for example, Tactical Resolution and Safety
Validation read the constraint zones of \fname{01b}.}
\label{fig:flow}
\end{figure}
